%% ODER: format ==         = "\mathrel{==}"
%% ODER: format /=         = "\neq "
%
%
\makeatletter
\@ifundefined{lhs2tex.lhs2tex.sty.read}%
  {\@namedef{lhs2tex.lhs2tex.sty.read}{}%
   \newcommand\SkipToFmtEnd{}%
   \newcommand\EndFmtInput{}%
   \long\def\SkipToFmtEnd#1\EndFmtInput{}%
  }\SkipToFmtEnd

\newcommand\ReadOnlyOnce[1]{\@ifundefined{#1}{\@namedef{#1}{}}\SkipToFmtEnd}
\usepackage{amstext}
\usepackage{amssymb}
\usepackage{stmaryrd}
\DeclareFontFamily{OT1}{cmtex}{}
\DeclareFontShape{OT1}{cmtex}{m}{n}
  {<5><6><7><8>cmtex8
   <9>cmtex9
   <10><10.95><12><14.4><17.28><20.74><24.88>cmtex10}{}
\DeclareFontShape{OT1}{cmtex}{m}{it}
  {<-> ssub * cmtt/m/it}{}
\newcommand{\texfamily}{\fontfamily{cmtex}\selectfont}
\DeclareFontShape{OT1}{cmtt}{bx}{n}
  {<5><6><7><8>cmtt8
   <9>cmbtt9
   <10><10.95><12><14.4><17.28><20.74><24.88>cmbtt10}{}
\DeclareFontShape{OT1}{cmtex}{bx}{n}
  {<-> ssub * cmtt/bx/n}{}
\newcommand{\tex}[1]{\text{\texfamily#1}}	% NEU

\newcommand{\Sp}{\hskip.33334em\relax}


\newcommand{\Conid}[1]{\mathit{#1}}
\newcommand{\Varid}[1]{\mathit{#1}}
\newcommand{\anonymous}{\kern0.06em \vbox{\hrule\@width.5em}}
\newcommand{\plus}{\mathbin{+\!\!\!+}}
\newcommand{\bind}{\mathbin{>\!\!\!>\mkern-6.7mu=}}
\newcommand{\rbind}{\mathbin{=\mkern-6.7mu<\!\!\!<}}% suggested by Neil Mitchell
\newcommand{\sequ}{\mathbin{>\!\!\!>}}
\renewcommand{\leq}{\leqslant}
\renewcommand{\geq}{\geqslant}
\usepackage{polytable}

%mathindent has to be defined
\@ifundefined{mathindent}%
  {\newdimen\mathindent\mathindent\leftmargini}%
  {}%

\def\resethooks{%
  \global\let\SaveRestoreHook\empty
  \global\let\ColumnHook\empty}
\newcommand*{\savecolumns}[1][default]%
  {\g@addto@macro\SaveRestoreHook{\savecolumns[#1]}}
\newcommand*{\restorecolumns}[1][default]%
  {\g@addto@macro\SaveRestoreHook{\restorecolumns[#1]}}
\newcommand*{\aligncolumn}[2]%
  {\g@addto@macro\ColumnHook{\column{#1}{#2}}}

\resethooks

\newcommand{\onelinecommentchars}{\quad-{}- }
\newcommand{\commentbeginchars}{\enskip\{-}
\newcommand{\commentendchars}{-\}\enskip}

\newcommand{\visiblecomments}{%
  \let\onelinecomment=\onelinecommentchars
  \let\commentbegin=\commentbeginchars
  \let\commentend=\commentendchars}

\newcommand{\invisiblecomments}{%
  \let\onelinecomment=\empty
  \let\commentbegin=\empty
  \let\commentend=\empty}

\visiblecomments

\newlength{\blanklineskip}
\setlength{\blanklineskip}{0.66084ex}

\newcommand{\hsindent}[1]{\quad}% default is fixed indentation
\let\hspre\empty
\let\hspost\empty
\newcommand{\NB}{\textbf{NB}}
\newcommand{\Todo}[1]{$\langle$\textbf{To do:}~#1$\rangle$}

\EndFmtInput
\makeatother
%
%
%
%
%
%
% This package provides two environments suitable to take the place
% of hscode, called "plainhscode" and "arrayhscode". 
%
% The plain environment surrounds each code block by vertical space,
% and it uses \abovedisplayskip and \belowdisplayskip to get spacing
% similar to formulas. Note that if these dimensions are changed,
% the spacing around displayed math formulas changes as well.
% All code is indented using \leftskip.
%
% Changed 19.08.2004 to reflect changes in colorcode. Should work with
% CodeGroup.sty.
%
\ReadOnlyOnce{polycode.fmt}%
\makeatletter

\newcommand{\hsnewpar}[1]%
  {{\parskip=0pt\parindent=0pt\par\vskip #1\noindent}}

% can be used, for instance, to redefine the code size, by setting the
% command to \small or something alike
\newcommand{\hscodestyle}{}

% The command \sethscode can be used to switch the code formatting
% behaviour by mapping the hscode environment in the subst directive
% to a new LaTeX environment.

\newcommand{\sethscode}[1]%
  {\expandafter\let\expandafter\hscode\csname #1\endcsname
   \expandafter\let\expandafter\endhscode\csname end#1\endcsname}

% "compatibility" mode restores the non-polycode.fmt layout.

\newenvironment{compathscode}%
  {\par\noindent
   \advance\leftskip\mathindent
   \hscodestyle
   \let\\=\@normalcr
   \let\hspre\(\let\hspost\)%
   \pboxed}%
  {\endpboxed\)%
   \par\noindent
   \ignorespacesafterend}

\newcommand{\compaths}{\sethscode{compathscode}}

% "plain" mode is the proposed default.
% It should now work with \centering.
% This required some changes. The old version
% is still available for reference as oldplainhscode.

\newenvironment{plainhscode}%
  {\hsnewpar\abovedisplayskip
   \advance\leftskip\mathindent
   \hscodestyle
   \let\hspre\(\let\hspost\)%
   \pboxed}%
  {\endpboxed%
   \hsnewpar\belowdisplayskip
   \ignorespacesafterend}

\newenvironment{oldplainhscode}%
  {\hsnewpar\abovedisplayskip
   \advance\leftskip\mathindent
   \hscodestyle
   \let\\=\@normalcr
   \(\pboxed}%
  {\endpboxed\)%
   \hsnewpar\belowdisplayskip
   \ignorespacesafterend}

% Here, we make plainhscode the default environment.

\newcommand{\plainhs}{\sethscode{plainhscode}}
\newcommand{\oldplainhs}{\sethscode{oldplainhscode}}
\plainhs

% The arrayhscode is like plain, but makes use of polytable's
% parray environment which disallows page breaks in code blocks.

\newenvironment{arrayhscode}%
  {\hsnewpar\abovedisplayskip
   \advance\leftskip\mathindent
   \hscodestyle
   \let\\=\@normalcr
   \(\parray}%
  {\endparray\)%
   \hsnewpar\belowdisplayskip
   \ignorespacesafterend}

\newcommand{\arrayhs}{\sethscode{arrayhscode}}

% The mathhscode environment also makes use of polytable's parray 
% environment. It is supposed to be used only inside math mode 
% (I used it to typeset the type rules in my thesis).

\newenvironment{mathhscode}%
  {\parray}{\endparray}

\newcommand{\mathhs}{\sethscode{mathhscode}}

% texths is similar to mathhs, but works in text mode.

\newenvironment{texthscode}%
  {\(\parray}{\endparray\)}

\newcommand{\texths}{\sethscode{texthscode}}

% The framed environment places code in a framed box.

\def\codeframewidth{\arrayrulewidth}
\RequirePackage{calc}

\newenvironment{framedhscode}%
  {\parskip=\abovedisplayskip\par\noindent
   \hscodestyle
   \arrayrulewidth=\codeframewidth
   \tabular{@{}|p{\linewidth-2\arraycolsep-2\arrayrulewidth-2pt}|@{}}%
   \hline\framedhslinecorrect\\{-1.5ex}%
   \let\endoflinesave=\\
   \let\\=\@normalcr
   \(\pboxed}%
  {\endpboxed\)%
   \framedhslinecorrect\endoflinesave{.5ex}\hline
   \endtabular
   \parskip=\belowdisplayskip\par\noindent
   \ignorespacesafterend}

\newcommand{\framedhslinecorrect}[2]%
  {#1[#2]}

\newcommand{\framedhs}{\sethscode{framedhscode}}

% The inlinehscode environment is an experimental environment
% that can be used to typeset displayed code inline.

\newenvironment{inlinehscode}%
  {\(\def\column##1##2{}%
   \let\>\undefined\let\<\undefined\let\\\undefined
   \newcommand\>[1][]{}\newcommand\<[1][]{}\newcommand\\[1][]{}%
   \def\fromto##1##2##3{##3}%
   \def\nextline{}}{\) }%

\newcommand{\inlinehs}{\sethscode{inlinehscode}}

% The joincode environment is a separate environment that
% can be used to surround and thereby connect multiple code
% blocks.

\newenvironment{joincode}%
  {\let\orighscode=\hscode
   \let\origendhscode=\endhscode
   \def\endhscode{\def\hscode{\endgroup\def\@currenvir{hscode}\\}\begingroup}
   %\let\SaveRestoreHook=\empty
   %\let\ColumnHook=\empty
   %\let\resethooks=\empty
   \orighscode\def\hscode{\endgroup\def\@currenvir{hscode}}}%
  {\origendhscode
   \global\let\hscode=\orighscode
   \global\let\endhscode=\origendhscode}%

\makeatother
\EndFmtInput
%
%% ----------------------
%% Stuff for TimCore


%% 'g' inferences




%% --------- Effects ------------


%% ----------------------




%% ----------- Applications --------------
%% %format applyx a b = a "@" b
%% %format applyx a b = a "\," b

%% Function application with no arguments

%% Function application with 1 argument

%% Function application with 1 argument, without parenthesis

%% Function application with many arguments

%% ----------- End of Applications --------------


%% Sequences e1; ...; en; v

%% xs etc stole syntax used in examples.

%% Old version: | for BNF, [] for choice
%%%format choice = "[ \! ]"
%%%format `choice` = "\mathop{[ \! ]}"
%% 1mm is about 0.3em with the current font

%% New version: tall | for BNF, fat | for choice
%%format vbar = "\mathop{\bigm|}"
%%format choice = "\boldsymbol{|}"
%%format `choice` = "\mathop{\boldsymbol{|}}"




%% Arrays and tuples
%% %format arrKW = "\textbf{arr}"
%% %format arr (x) = arrKW "\{" x "\}"





%%  OLD SYNTAX   %format split (x) (y) (z) = "\textbf{split}(" x ")\{" y "," z "\}"
%% %format defKW = "\textbf{def}"

% only used in comments, but still prettier this way:
%%%%format map = "\textbf{map}"

%% Effects checking, like succeeds{e}


%% Unification
%%%format == = "\!=\!"


% format := = "\mapsto"
%% %format squiggt = "\!\leadsto\!"
%  format effs x = "\!\!<\!\!" x "\!\!>\!\!" dots


































%% ----------- Metatheory --------------
%%Formatting for skew confluence proof


%%%%%%%%%%%%%%%%%%%%%%%%%%%%%%%%%%%%

\section{Mathematical Background}


Our development of the background material on relations, terms, substitutions, and rewrite rules substantially follows that of
\citet{Huet80}, with considerable influence also from 
\citet{van-oostrom:decreasing-diagrams} and \citet{ariola:skew-confluence}.


\subsection{Relations, Terms, and Rewrite Rules}


\subsubsection{Relations}


A \emph{relation} is a set of pairs of related items. We use letters, such as $R$ and $S$, as well as symbols, such as $\subset$ and $\leq$ and $\rightarrow$ and $\equiv$, to name relations.
If $A$ and $B$ are sets, then any $R \subseteq A \times B$ is a \emph{relation from $A$ to $B$}; if $A=B$, so that $R \subseteq A \times A$,
then $R$ is a \emph{relation on $A$}.

Any relation, regarded as a set, may be used as the operand of operators such as $\cup$ and $\cap$ and of predicates such as $\subseteq$ and $\in$.
Any relation may also be regarded as a predicate; if $R$ is a relation, then $x \mathrel{R} y$ means $(x,y)\in R$.
Because letters and symbols alike may be used to name relations, we can also say that if $\rightarrow$ is a relation, then $x \mathrel{\rightarrow} y$ means $(x, y)\in {\rightarrow}$.
We also use the customary notation for \emph{chaining} relation predicates, so that for any $n>0$, $x_0 \mathrel{R_1} x_1 \mathrel{R_2}  x_2 \mathrel{R_3} \cdots \mathrel{R_n} x_n$
means $x_0 \mathrel{R_1} x_1 \logand x_1 \mathrel{R_2} x_2 \logand x_2 \mathrel{R_3} x_3 \logand \cdots \logand x_{n-1} \mathrel{R_n} x_n$.
For example, $x \leq y \leq z$ means $x \leq y \logand y \leq z$.


We write $\iota_A$ for the identity relation on $A$: $\iota_A = \setcomp{(x,x)}{x \in A}$.
The operator $\relcomp$ is \emph{relation composition}: if $R$ and $S$ are relations, then $R \relcomp S$ means $\setcomp{(x, y)}{\exists z \suchthat x \mathrel{R} z \mathrel{S} y}$.
A nonnegative integer superscript $k$ may used to indicate $k$-fold composition of a relation; if $R$ and $\xgo{}$ are relations on $A$, then $R^0 = {\xgosup{}{0}} = \iota_A$,
and for $k>0$, $R^{k} = R \relcomp R^{k-1}$ and ${\xgosup{}{k}} = {\xgo{}} \relcomp {\xgosup{}{k-1}}$.


The \emph{inverse} of a relation $R$ is the set $R^{-1} = \setcomp{(x,y)}{(y,x) \in R}$; note that the inverse of the inverse of $R$ is $R$.
While the superscript $^{-1}$ could be applied to symbols as well as letters, we instead customarily notate the inverse of a relation
that is named by a symbol by flipping the symbol left-to-right, so that $\supset$ denotes the inverse of $\subset$, $\geq$ denotes the inverse of $\leq$,
and $\leftarrow$ denotes the inverse of $\rightarrow$.

A relation $R$ is said to be \emph{symmetric} if $R^{-1}=R$. It is customary to use symbols that are left-right symmetric, such as $\equiv$, to denote symmetric relations.

Again, let $R$ and $\xgo{}$ be relations on $A$.
The reflexive closure of $R$ is $R^{\reflexivemarker}=R \cup \iota_A$ and the reflexive closure of $\xgo{}$ is ${\xgoeq{}} = {\xgo{}} \cup \iota_A$.
The transitive reflexive closure of $R$ is $R^{*}= \bigcup_{k \geq 0} R^k$ and the transitive reflexive closure of $\xgo{}$ is ${\xgodbl{}} = \bigcup_{k \geq 0} \xgosup{}{k}$.
The symmetric closure of $\xgo{}$ is ${\xboth{}} = {\xungo{}} \cup {\xgo{}}$; 
the symmetric transitive reflexive closure of $\xgo{}$ is ${\xbothdbl{}} = {\xungodbl{}} \cup {\xgodbl{}}$. 

For a given relation denoted by a rightward arrow $\xgo{}$, we define the predicate $x \bothreduceto y$ to mean $\exists z \suchthat x \xgodbl{} z \xungodbl{} y$
and the predicate $x \bothreducefrom y$ to mean $\exists z \suchthat x \xungodbl{} z \xgodbl{} y$.

\subsubsection{Terms and Occurrences}

Our definition of \emph{term} is specialized for the grammar of {\versecalc}. Briefly put, a term is the parse tree for any BNF nonterminal in the {\versecalc} grammar
or any string into which such a nonterminal might be expanded, either partially or fully. For our purposes here, we require that terms be finite.
Let {\metavars} be the set of of \emph{metavariables} (possibly decorated BNF nonterminals) and
let {\termset} be the set of all possible terms.
For any term $M \in \termset$, we define its \emph{set of metavariables} $\metavars(M) \subset \metavars$, as shown in \cref{fig:metavars}.

We use elements of $\mathbb{N}_{+}^{*}$ (that is, tuples of positive integers) to represent the positions of the subterms of a term.
We call such tuples \emph{occurrences}; we use variables $u$, $w$, $p$, and $q$ to refer to them,
and we use variables $U$, $W$, $P$, and $Q$ to refer to sets of them.
We use juxtaposition to concatenate such tuples; for example, if $u=(1,2)$ and $w=(2,1,3)$, then $uw =(1,2,2,1,3)$ and $w(2,2)u = (2,1,3,2,2,1,2)$.
Furthermore we sometimes use juxtaposition to concatenate a sequence with every member of a set to form a new set;
for example, if $U=\{(1,2),(2,1,3)\}$, then $(2)U = \{(2,1,2),(2,2,1,3)\}$.
We say \emph{$u$ is within $w$}, written $u \occin w$, iff $\exists q \in \mathbb{N}^{*}_{+} \suchthat u=wq$,
in which case we also define $u/w = q$.\footnote{This overloading of the operator $/$ allows
  a very pretty definition of the distributivity property for replacement in \cref{sec:replacement}.}
We say \emph{$u$ and $w$ are disjoint}, written $u \occdis w$, iff $\neg\bigl( u \occin w \logor w \occin u \bigr)$.
Finally, $u \occstrictin w$ means $u \occin w \logand u \neq w$ and $u \occdiseq w$ means $u \occdis w \logor u = w$.

For any term $M \in \termset$, we define its \emph{set of occurrences} $\occ(M) \subset \mathbb{N}_{+}^{*}$, as shown in \cref{fig:occurrences}, and the \emph{subterm of $M$ at $u$},
$M/u \in \termset$, where $u \in \occ(M)$, as shown in \cref{fig:selection}. We overload the relations $\occin$ and $\occstrictin$ to apply to terms:
$M \occin N$ means $\exists u \suchthat M = N/u$, 
and $M \occstrictin N$ means $\exists u \neq \emptyocc \suchthat M = N/u$.
We also say that \emph{$u$ is an occurrence of the term $M/u$ in $M$}.

We define $\occ(M,N)$ (the \emph{occurrences of $M$ in $N$} to mean $\setcomp{u}{u\in\occ(N) \logand N/u=M}$.  
We define $\uniqueocc(M,N)$ (the \emph{unique occurrence of $M$ in $N$} to mean $u$ where $\occ(M,N) = \{u\}$;
$\uniqueocc(M,N)$ is undefined if $\left|\occ(M,N)\right| \neq 1$.


\begin{figure}
$
\begin{array}{@{}l@{\qquad}l@{\qquad}l@{}}
  \hbox{form of $M$} & \metavars(M) & \occ(M) \\[4pt] \hline
  M \in \metavars & \{M\} & \{\emptyocc\} \\ \hline
  \ensuremath{\Varid{x}} & \{\,\} & \{\emptyocc\} \\ \hline
  \ensuremath{\Varid{k}} & \{\,\} & \{\emptyocc\} \\ \hline
  \ensuremath{\Varid{op}} & \{\,\} & \{\emptyocc\} \\ \hline
  \ensuremath{\langle{}\Varid{v}_{\mathrm{1}},\;\ldots,\Varid{v}_{\Varid{n}}\rangle} & \bigcup_{1\leq i\leq n} \metavars(\ensuremath{\Varid{v}_{\Varid{i}}})& \{\emptyocc\} \cup \bigcup_{1\leq i\leq n} (i)\occ(\ensuremath{\Varid{v}_{\Varid{i}}}) \\ \hline
  \ensuremath{\lambda\Varid{x}.\,\Varid{e}} & \metavars(\ensuremath{\Varid{x}}) \cup \metavars(\ensuremath{\Varid{e}}) & \{\emptyocc\} \cup (1)\occ(\ensuremath{\Varid{x}}) \cup (2)\occ(\ensuremath{\Varid{e}}) \\ \hline
  \ensuremath{\mu\Varid{x}.\,\Varid{e}} & \metavars(\ensuremath{\Varid{x}}) \cup \metavars(\ensuremath{\Varid{e}}) & \{\emptyocc\} \cup (1)\occ(\ensuremath{\Varid{x}}) \cup (2)\occ(\ensuremath{\Varid{e}}) \\ \hline
  \ensuremath{eq;\,\Varid{e}} & \metavars(\ensuremath{eq}) \cup \metavars(\ensuremath{\Varid{e}}) & \{\emptyocc\} \cup (1)\occ(\ensuremath{eq}) \cup (2)\occ(\ensuremath{\Varid{e}}) \\ \hline
  \ensuremath{\Varid{v}\mathrel{=}\Varid{e}} & \metavars(\ensuremath{\Varid{v}}) \cup \metavars(\ensuremath{\Varid{e}}) & \{\emptyocc\} \cup (1)\occ(\ensuremath{\Varid{v}}) \cup (2)\occ(\ensuremath{\Varid{e}}) \\ \hline
  \ensuremath{\exists\Varid{x}.\,\Varid{e}} & \metavars(\ensuremath{\Varid{x}}) \cup \metavars(\ensuremath{\Varid{e}}) & \{\emptyocc\} \cup (1)\occ(\ensuremath{\Varid{x}}) \cup (2)\occ(\ensuremath{\Varid{e}}) \\ \hline
  \ensuremath{\textbf{fail}} & \{\,\} & \{\emptyocc\} \\ \hline
  \ensuremath{\Varid{e}_{\mathrm{1}}\;\rule[-0.1em]{0.15em}{0.75em}\;\Varid{e}_{\mathrm{2}}} & \metavars(\ensuremath{\Varid{e}_{\mathrm{1}}}) \cup \metavars(\ensuremath{\Varid{e}_{\mathrm{2}}}) & \{\emptyocc\} \cup (1)\occ(\ensuremath{\Varid{e}_{\mathrm{1}}}) \cup (2)\occ(\ensuremath{\Varid{e}_{\mathrm{2}}}) \\ \hline
  \ensuremath{\Varid{v}_{\mathrm{1}}\,\Varid{v}_{\mathrm{2}}} & \metavars(\ensuremath{\Varid{v}_{\mathrm{1}}}) \cup \metavars(\ensuremath{\Varid{v}_{\mathrm{2}}}) & \{\emptyocc\} \cup (1)\occ(\ensuremath{\Varid{v}_{\mathrm{1}}}) \cup (2)\occ(\ensuremath{\Varid{v}_{\mathrm{2}}}) \\ \hline
  \ensuremath{\textbf{one}\{\Varid{e}\}} & \metavars(\ensuremath{\Varid{e}}) & \{\emptyocc\} \cup (1)\occ(\ensuremath{\Varid{e}}) \\ \hline
  \ensuremath{\textbf{all}\{\Varid{e}\}} & \metavars(\ensuremath{\Varid{e}}) & \{\emptyocc\} \cup (1)\occ(\ensuremath{\Varid{e}}) \\ \hline
\end{array}
$
\caption{Definitions of $\metavars(M)$ (set of metavariables within $M$) and $\occ(M)$ (the set of occurrences within $M$)}\label{fig:metavars}\label{fig:occurrences}
\Description{Definition of notation for the set of occurrences of subterms within a term $M$}
\end{figure}

\begin{figure}
$
\begin{array}{@{}l@{\qquad}l@{\quad}l@{}}
  \hbox{form of $M$} & \hbox{meaning of $M/u$} \\[4pt] \hline
  \hbox{any} & M/\emptyocc = M \\ \hline
  \ensuremath{\langle{}\Varid{v}_{\mathrm{1}},\;\ldots,\Varid{v}_{\Varid{n}}\rangle} & M/(i)u = \ensuremath{\Varid{v}_{\Varid{i}}}/u & \hbox{for all $1\leq i\leq n, u\in\occ(\ensuremath{\Varid{v}_{\Varid{i}}})$} \\ \hline
  \ensuremath{\lambda\Varid{x}.\,\Varid{e}} & M/(1)u = \ensuremath{\Varid{x}}/u & \hbox{for all $u\in\occ(\ensuremath{\Varid{x}})$} \\ & M/(2)u = \ensuremath{\Varid{e}}/u & \hbox{for all $u\in\occ(\ensuremath{\Varid{e}})$} \\ \hline
  \ensuremath{\mu\Varid{x}.\,\Varid{e}} & M/(1)u = \ensuremath{\Varid{x}}/u & \hbox{for all $u\in\occ(\ensuremath{\Varid{x}})$} \\ & M/(2)u = \ensuremath{\Varid{e}}/u & \hbox{for all $u\in\occ(\ensuremath{\Varid{e}})$} \\ \hline
  \ensuremath{eq;\,\Varid{e}} & M/(1)u = \ensuremath{eq}/u & \hbox{for all $u\in\occ(\ensuremath{eq})$} \\ & M/(2)u = \ensuremath{\Varid{e}}/u & \hbox{for all $u\in\occ(\ensuremath{\Varid{e}})$} \\ \hline
  \ensuremath{\Varid{v}\mathrel{=}\Varid{e}} & M/(1)u = \ensuremath{\Varid{v}}/u & \hbox{for all $u\in\occ(\ensuremath{\Varid{v}})$} \\ & M/(2)u = \ensuremath{\Varid{e}}/u & \hbox{for all $u\in\occ(\ensuremath{\Varid{e}})$} \\ \hline
  \ensuremath{\exists\Varid{x}.\,\Varid{e}} & M/(1)u = \ensuremath{\Varid{x}}/u & \hbox{for all $u\in\occ(\ensuremath{\Varid{x}})$} \\ & M/(2)u = \ensuremath{\Varid{e}}/u & \hbox{for all $u\in\occ(\ensuremath{\Varid{e}})$} \\ \hline
  \ensuremath{\Varid{e}_{\mathrm{1}}\;\rule[-0.1em]{0.15em}{0.75em}\;\Varid{e}_{\mathrm{2}}} & M/(1)u = \ensuremath{\Varid{e}_{\mathrm{1}}}/u & \hbox{for all $u\in\occ(\ensuremath{\Varid{e}_{\mathrm{1}}})$} \\ & M/(2)u = \ensuremath{\Varid{e}_{\mathrm{2}}}/u & \hbox{for all $u\in\occ(\ensuremath{\Varid{e}_{\mathrm{2}}})$} \\ \hline
  \ensuremath{\Varid{v}_{\mathrm{1}}\,\Varid{v}_{\mathrm{2}}} & M/(1)u = \ensuremath{\Varid{v}_{\mathrm{1}}}/u & \hbox{for all $u\in\occ(\ensuremath{\Varid{v}_{\mathrm{1}}})$} \\ & M/(2)u = \ensuremath{\Varid{v}_{\mathrm{2}}}/u & \hbox{for all $u\in\occ(\ensuremath{\Varid{v}_{\mathrm{2}}})$} \\ \hline
  \ensuremath{\textbf{one}\{\Varid{e}\}} & M/(1)u = \ensuremath{\Varid{e}}/u & \hbox{for all $u\in\occ(\ensuremath{\Varid{e}})$} \\ \hline
  \ensuremath{\textbf{all}\{\Varid{e}\}} & M/(1)u = \ensuremath{\Varid{e}}/u & \hbox{for all $u\in\occ(\ensuremath{\Varid{e}})$} \\ \hline
\end{array}
$
\caption{Definition of $M/u$ for a term $M$ and $u \in \occ(M)$}\label{fig:selection}
\Description{Definitions of notation for selection of a subterm at occurrence $u$ within a term $M$}
\end{figure}

\newpage

\subsubsection{Replacement of Subterms}\label{sec:replacement}


For any $M\in\termset$, $u \in \occ(M)$, and $N \in\termset$, we define the \emph{replacement} operation \ensuremath{\tsubst{\Conid{M}}{\Varid{u}}{\Conid{N}}} as shown in \cref{fig:replacement};
note that $N$ is required to be a term of the correct ``type,'' so that the result of  \ensuremath{\tsubst{\Conid{M}}{\Varid{u}}{\Conid{N}}} will be a member of $\termset$.
The result is a term that is a copy of $M$ except that the subterm at $u$ is replaced by $N$.

Replacement has the following properties~\cite[Proposition 3.1]{Huet80}:
\begin{quote}
$\forall M,N,P \in\termset, u \in \occ(M), w \in \occ(N)$:
\begin{tabbing}
\qquad\=\hskip 3in\=\kill 
\>$\tsubst{M}{u}{N}/uw = N/w$ \>{\rm\emph{embedding}} \\  
\>$\tsubst{\tsubst{M}{u}{N}}{uw}{P} = \tsubst{M}{u}{\tsubst{N}{w}{P}}$ \>{\rm\emph{associativity}}
\end{tabbing}
$\forall M,N,P \in\termset, u,w \in \occ(M)$ such that $u \occdis w$:
\begin{tabbing}
\qquad\=\hskip 3in\=\kill 
\>$\tsubst{M}{u}{N}/w = M/w$ \>{\rm\emph{persistence}} \\  
\>$\tsubst{\tsubst{M}{u}{N}}{w}{P} = \tsubst{\tsubst{M}{w}{P}}{u}{N}$ \>{\rm\emph{commutativity}}
\end{tabbing}
$\forall M,N,P \in\termset, u,w \in \occ(M)$ such that $u \occin w$:
\begin{tabbing}
\qquad\=\hskip 3in\=\kill 
\>$\tsubst{M}{u}{N}/w = \tsubst{(M/w)}{u/w}{N}$ \>{\rm\emph{distributivity}} \\  
\>$\tsubst{\tsubst{M}{u}{N}}{w}{P} = \tsubst{M}{w}{P}$ \>{\rm\emph{dominance}}
\end{tabbing}
\end{quote}

We also propose names for two additional properties:

\begin{lemma}
$\forall M,N,P \in\termset, u \in \occ(M), w \in \occ(N)$:
\begin{quote}
\begin{tabbing}
\qquad\=\hskip 3in\=\kill 
\>$\tsubst{M}{u}{M/u} = M$\>{\rm\emph{consistency}}
\end{tabbing}
\end{quote}
\end{lemma}

\begin{lemma}
$\forall M,N \in\termset, q,u \in \occ(M)$ such that $q \occin u$, and let $w=q/u$:\hfill\break
\begin{quote}
\begin{tabbing}
\qquad\=\hskip 3in\=\kill 
\>$\tsubst{M}{uw}{N} = \tsubst{M}{u}{\tsubst{(M/u)}{w}{N}}$\>{\rm\emph{factoring}}
\end{tabbing}
\end{quote}
\end{lemma}
\begin{proof}
  $\tsubst{M}{uw}{N} = \tsubst{\tsubst{M}{u}{M/u}}{uw}{N} \BY{\emph{consistency}}$ \hfill\break \null\hskip1.35in${} = \tsubst{M}{u}{\tsubst{(M/u)}{w}{N}} \BY{\emph{associativity}}$.
\end{proof}




For any $M\in\termset$, any $U \subseteq \occ(M)$ such that $\forall u, w \in U \suchthat \bigl( M/u = M/w \logand u \occdiseq w\bigr)$, and any $N \in\termset$,
we define \ensuremath{\tsubst{\Conid{M}}{\Conid{U}}{\Conid{N}}}
to mean $(\cdots\tsubst{(\tsubst{(\tsubst{M}{u_1}{N})}{u_2}{N})\cdots)}{u_n}{N}$ for any arbitrarily chosen ordering $u_1, u_2, \ldots, u_n$ of the elements of $U$.
The result is a copy of $M$ with any number of its subterms replaced by copies of $N$, but the subterms replaced must all be the same and their occurrences must be disjoint within $M$
(these restrictions plus the commutativity property for replacement imply that the result is the same no matter what ordering is chosen for the elements of $U$).


\begin{figure}
$
\begin{array}{@{}l@{\hskip1.5em}l@{\hskip1.5em}l@{}}
  \hbox{form of $M$} & \hbox{meaning of \ensuremath{\tsubst{\Conid{M}}{\Varid{u}}{\Conid{N}}}} \\[4pt] \hline
  \hbox{any} & \ensuremath{\tsubst{\Conid{M}}{()}{\Conid{N}}} \;\;=\;\; N \\ \hline
  \ensuremath{\langle{}\Varid{v}_{\mathrm{1}},\;\ldots,\Varid{v}_{\Varid{n}}\rangle} & \tsubst{M}{(i)u}{N} \;\;=\;\; \ensuremath{\langle{}\Varid{v}_{\mathrm{1}},\;\ldots,v_{i-1},\tsubst{\Varid{v}_{\Varid{i}}}{\Varid{u}}{\Conid{N}},v_{i+1},\;\ldots,\Varid{v}_{\Varid{n}}\rangle} & \hbox{for all $1\leq i\leq n$} \\ \hline
  \ensuremath{\lambda\Varid{x}.\,\Varid{e}}  & \tsubst{M}{(1)u}{N} \;\;=\;\; \ensuremath{\lambda\tsubst{\Varid{x}}{\Varid{u}}{\Conid{N}}.\,\Varid{e}} \\ & \tsubst{M}{(2)u}{N} \;\;=\;\; \ensuremath{\lambda\Varid{x}.\,\tsubst{\Varid{e}}{\Varid{u}}{\Conid{N}}} \\ \hline
  \ensuremath{\mu\Varid{x}.\,\Varid{e}}  & \tsubst{M}{(1)u}{N} \;\;=\;\; \ensuremath{\mu\tsubst{\Varid{x}}{\Varid{u}}{\Conid{N}}.\,\Varid{e}} \\ & \tsubst{M}{(2)u}{N} \;\;=\;\; \ensuremath{\mu\Varid{x}.\,\tsubst{\Varid{e}}{\Varid{u}}{\Conid{N}}} \\ \hline
  \ensuremath{eq;\,\Varid{e}}  & \tsubst{M}{(1)u}{N} \;\;=\;\; \ensuremath{\tsubst{eq}{\Varid{u}}{\Conid{N}};\,\Varid{e}} \\ & \tsubst{M}{(2)u}{N} \;\;=\;\; \ensuremath{eq;\,\tsubst{\Varid{e}}{\Varid{u}}{\Conid{N}}} \\ \hline
  \ensuremath{\Varid{v}\mathrel{=}\Varid{e}}  & \tsubst{M}{(1)u}{N} \;\;=\;\; \ensuremath{\tsubst{\Varid{v}}{\Varid{u}}{\Conid{N}}\mathrel{=}\Varid{e}} \\ & \tsubst{M}{(2)u}{N} \;\;=\;\; \ensuremath{\Varid{v}\mathrel{=}\tsubst{\Varid{e}}{\Varid{u}}{\Conid{N}}} \\ \hline
  \ensuremath{\exists\Varid{x}.\,\Varid{e}}  & \tsubst{M}{(1)u}{N} \;\;=\;\; \ensuremath{\exists\tsubst{\Varid{x}}{\Varid{u}}{\Conid{N}}.\,\Varid{e}} \\ & \tsubst{M}{(2)u}{N} \;\;=\;\; \ensuremath{\exists\Varid{x}.\,\tsubst{\Varid{e}}{\Varid{u}}{\Conid{N}}} \\ \hline
  \ensuremath{\Varid{e}_{\mathrm{1}}\;\rule[-0.1em]{0.15em}{0.75em}\;\Varid{e}_{\mathrm{2}}}  & \tsubst{M}{(1)u}{N} \;\;=\;\; \ensuremath{\tsubst{\Varid{e}_{\mathrm{1}}}{\Varid{u}}{\Conid{N}}\;\rule[-0.1em]{0.15em}{0.75em}\;\Varid{e}_{\mathrm{2}}} \\ & \tsubst{M}{(2)u}{N} \;\;=\;\; \ensuremath{\Varid{e}_{\mathrm{1}}\;\rule[-0.1em]{0.15em}{0.75em}\;\tsubst{\Varid{e}_{\mathrm{2}}}{\Varid{u}}{\Conid{N}}} \\ \hline
  \ensuremath{\Varid{v}_{\mathrm{1}}\,\Varid{v}_{\mathrm{2}}}  & \tsubst{M}{(1)u}{N} \;\;=\;\; \ensuremath{\tsubst{\Varid{v}_{\mathrm{1}}}{\Varid{u}}{\Conid{N}}\,\Varid{v}_{\mathrm{2}}} \\ & \tsubst{M}{(2)u}{N} \;\;=\;\; \ensuremath{\Varid{v}_{\mathrm{1}}\,\tsubst{\Varid{v}_{\mathrm{2}}}{\Varid{u}}{\Conid{N}}} \\ \hline
  \ensuremath{\textbf{one}\{\Varid{e}\}}  & \tsubst{M}{(1)u}{N} \;\;=\;\; \ensuremath{\textbf{one}\{\tsubst{\Varid{e}}{\Varid{u}}{\Conid{N}}\}} \\ \hline
  \ensuremath{\textbf{all}\{\Varid{e}\}}  & \tsubst{M}{(1)u}{N} \;\;=\;\; \ensuremath{\textbf{all}\{\tsubst{\Varid{e}}{\Varid{u}}{\Conid{N}}\}} \\ \hline
\end{array}
$
\caption{Definition of $\tsubst{M}{u}{N}$ for a term $M$, $u \in \occ(M)$, and a term $N$}\label{fig:replacement}
\Description{Definitions of notation for replacing a subterm within a term $M$}
\end{figure}

\subsubsection{Substitutions for Metavariables}

\subsubsection{Matching of Terms}

\subsubsection{Rewrite Rules on Terms}

A \emph{rewrite rule} is a pair of terms, customarily written $\langle\alpha \xgo{} \beta\rangle$, such that $\metavars(\beta) \subseteq \metavars(\alpha)$.
(When rewrite rules are listed in a table, the surrounding angle brackets are typically omitted.)
A rewrite rule may be given a name, which we write in small caps; we also use the names \rulename{r} and \rulename{s} as metavariables
that range over rule names. A rewrite rule describes a relation on terms: $(M,N)$ is a member of the relation iff
$\exists \sigma \suchthat M = \sigma(\alpha) \logand N=\sigma(\beta)$.

Let $\xgo{}$ be a relation on $\termset$. We say that $\xgo{}$ is \emph{stable} iff
for all substitutions $\sigma$ and all terms $M$ and $N$ in $\termset$, if $M \xgo{} N$ then $\sigma(M) \xgo{} \sigma(N)$ (that is, the relationship is unaffected
if metavariables are renamed). We say that $\xgo{}$ is \emph{compatible} iff
for all terms $M$, $N$, and $P$ in $\termset$ and all occurrences $u \in \occ(P)$, if $M \xgo{} N$ then $\tsubst{P}{u}{M} \xgo{} \tsubst{P}{u}{N}$ (that is, a subterm can be
``rewritten in place'' in the same way as an entire term).
For any rewrite rule \rulename{r}, we write $\xrngo{r}$ to denote the smallest stable, compatible relation that contains the relation described by \rulename{r}.

If $\rr$ is a set of names of rewrite rules, we write $\xgo{\rr}$ to mean $\bigcup_{\hbox{\rulename{r}}\in\rr} \xrngo{r}$.

When we use $\xrngo{r}$ or $\xgo{\rr}$ to describe a relationship between two terms $M$ and $N$,
we sometimes annotate it to indicate explicitly the occurrence $u$ within $M$ of the subterm that is rewritten: $M \xrngosup{r}{u} N$ means
$M \xrngo{r} N \logand \exists \sigma \suchthat M/u = \sigma(\alpha) \logand N = \tsubst{M}{u}{\sigma(\beta)}$.

\begin{lemma}\label{lem:rewrite-basic-tsubst}
For any rewrite rule $\langle\alpha \xgo{} \beta\rangle$ (call it \rulename{r}) and terms $M$ and $N$,
\begin{quote}
\begin{tabbing}
\qquad\=\hskip 3in\=\kill 
\>if\quad$M \xrngosup{r}{u} N$\quad then\quad$N = \tsubst{M}{u}{N/u}$\>{\rm\emph{rewriting}}
\end{tabbing}
\end{quote}
\end{lemma}
\begin{proof}
  By the definition of $M \xrngosup{r}{u} N$ we have $\exists \sigma \suchthat M/u = \sigma(\alpha) \logand N = \tsubst{M}{u}{\sigma(\beta)}$.
  Now $N/u = \tsubst{M}{u}{\sigma(\beta)}/u = \sigma(\beta) \BY{\emph{embedding} with $w=\emptyocc$}$;
  therefore $N = \tsubst{M}{u}{N/u}$.
\end{proof}

We write $\xrngoeq{r}$ for the reflexive closure of $\xrngo{r}$, and we sometimes annotate it as $\xrngoeqsup{r}{u}$.

We write $\xrngodbl{r}$ for the transitive reflexive closure of $\xrngo{r}$, and similarly write  $\xgodbl{\rr}$ for the transitive reflexive closure of $\xgo{\rr}$.
We do not annotate such relations with occurrences $u$.\footnote{One could contemplate annotating them
  with \emph{sequences} of occurrences, but we will not need such annotations for our purposes here.}

We notate the inverses of these relations as
$\xrnungo{r}$ and
$\xungo{\rr}$ and
$\xrnungoeq{r}$ and
$\xungoeq{\rr}$ and
$\xrnungodbl{r}$ and
$\xungodbl{\rr}$, and we sometimes use annotated forms
$\xrnungosup{r}{u}$ and
$\xungosup{\rr}{u}$ and
$\xrnungoeqsup{r}{u}$ and
$\xungoeqsup{\rr}{u}$.

We notate the symmetric closures of these relations as
$\xrnboth{r}$ and
$\xboth{\rr}$ and
$\xrnbotheq{r}$ and
$\xbotheq{\rr}$ and
$\xrnbothdbl{r}$ and
$\xbothdbl{\rr}$, and we sometimes use annotated forms
$\xrnbothsup{r}{u}$ and
$\xbothsup{\rr}{u}$ and
$\xrnbotheqsup{r}{u}$ and
$\xbotheqsup{\rr}{u}$.


\subsubsection{Disjoint Parallel Compatible Relations}

The \emph{disjoint parallel closure} of the stable compatible relation $\xrngo{r}$ defined by a rewrite rule \rulename{r} is written $\xprngo{r}$;
it is the smallest relation containing $\xrngo{r}$ such that
for all terms $M$ and $N$ and sets of occurrences $U\in\occ(M)$ such that $\forall u, w\in U \suchthat u \occdiseq w$,
if $M \xprngo{r} N$ then $\tsubst{P}{U}{M} \xprngo{r} \tsubst{P}{U}{N}$ (that is, an entire set of identical subterms can each be
``rewritten in place'' in the same way as an entire term). Because the set $U$ may be empty, $\xprngo{r}$ is reflexive.

If $\rr$ is a set of names of rewrite rules, we write $\xpgo{\rr}$ to mean $\bigcup_{\hbox{\rulename{r}}\in\rr} \xprngo{r}$.

When we use $\xprngo{r}$ to describe a relationship between two terms $M$ and $N$,
we sometimes annotate it to indicate explicitly the set $U$ of occurrences within $M$ of the subterms that are rewritten: $M \xprngosup{r}{U} N$ means
$M \xprngo{r} N \logand \exists \sigma \suchthat \bigl(\forall u\in U\suchthat M/u = \sigma(\alpha)\bigr) \logand N = \tsubst{M}{U}{\sigma(\beta)}$.

We write $\xprngodbl{r}$ for the transitive reflexive closure of $\xprngo{r}$, and similarly write  $\xpgodbl{\rr}$ for the transitive reflexive closure of $\xpgo{\rr}$.
We do not annotate such relations with sets of occurrences $U$.\footnote{One could contemplate annotating them
  with \emph{sequences} of sets of occurrences, but we will not need such annotations for our purposes here.}

We notate the inverses of these relations as
$\xprnungo{r}$ and
$\xpungo{\rr}$ and
$\xprnungodbl{r}$ and
$\xpungodbl{\rr}$, and we sometimes use annotated forms
$\xprnungosup{r}{U}$ and
$\xpungosup{\rr}{U}$.

We notate the symmetric closures of these relations as
$\xprnboth{r}$ and
$\xpboth{\rr}$ and
$\xprnbothdbl{r}$ and
$\xpbothdbl{\rr}$, and we sometimes use annotated forms
$\xprnbothsup{r}{U}$ and
$\xpbothsup{\rr}{U}$.


\subsubsection{Free Variables}


We use the conventional notation $\freevars{\ensuremath{\Varid{e}}}$ to denote the set of variables that occur free
in the expression \ensuremath{\Varid{e}}. Variables are bound by the constructs \ensuremath{\exists\Varid{x}.\,\Varid{e}}, \ensuremath{\lambda\Varid{x}.\,\Varid{e}}, and \ensuremath{\mu\Varid{x}.\,\mathit{hnf}}.
\Cref{fig:fvs} contains a formal definition of this function for \versecalc{}.

We define $\freeocc(x,M)$ (the \emph{free occurrences of variable $x$ within $M$} as shown in \cref{fig:freeocc}.


We use the variation $\freevarsoutsidelambdas{\ensuremath{\Varid{e}}}$ to denote the set of variables that occur free
in the expression \ensuremath{\Varid{e}} in at least one position that is not within the body of a
lambda expression\footnote{
  ``$\freevarsoutsidelambdas{\cdot}$'' abbreviates ``free variables outside lambda''}.
As an example, $\freevarsoutsidelambdas{\ensuremath{\exists\Varid{x}.\,\langle{}\Varid{x},\Varid{f},\Varid{g},\lambda\Varid{y}.\,\langle{}\Varid{x},\Varid{g},\Varid{y},\Varid{z}\rangle\rangle}} = \{ f, g \}$ because:
\begin{itemize}
\item \ensuremath{\Varid{x}} is bound by \ensuremath{\exists{}}, so it is not free. 
\item \ensuremath{\Varid{f}} is free in a position not within the body of a lambda expression.
\item \ensuremath{\Varid{g}} is free in at least one position not within the body of a lambda expression (it also happens to be free in a second position that is within the body of a lambda expression).
\item \ensuremath{\Varid{y}} is bound by \ensuremath{\lambda\!}, so it is not free. 
\item \ensuremath{\Varid{z}} is free, but appears only in a position that is within the body of a lambda expression.
\end{itemize}
\Cref{fig:fvsol} contains a formal definition of this function.
Unlike $\freevars{e}$, $\freevarsoutsidelambdas{\ensuremath{\Varid{v}}}$ is only ever appplied to a value $\ensuremath{\Varid{v}}$.

\subsubsection{Substitution}

We use the notation \ensuremath{\subst{\Varid{e}}{\Varid{x}}{\Varid{v}}} to denote the expression that results from performing standard capture-avoiding substitution of the value \ensuremath{\Varid{v}} for every occurrence of the variable \ensuremath{\Varid{x}} within the expression \ensuremath{\Varid{e}} (it turns out that, for \versecalc{}, substitution of general expressions for variables is not required, only substitution of values for variables).
\Cref{fig:substitution} contains a formal definition of how this notation applies to the \versecalc{} grammar (compare \cite[\S 2.1.15]{Barendregt}).



\subsection{Skew Confluence}


\begin{comment}
\subsubsection{Overview of Skew Confluence (VERSION 1)} \label{sec:skew-overview}


If a reduction relation $\rr$ is confluent, then for all \ensuremath{\Varid{e}}, \ensuremath{\Varid{e}_{\mathrm{1}}}, and \ensuremath{\Varid{e}_{\mathrm{2}}},
if $\rstep{\rr}{\ensuremath{\Varid{e}}}{\ensuremath{\Varid{e}_{\mathrm{1}}}}$ and $\rstep{\rr}{\ensuremath{\Varid{e}}}{\ensuremath{\Varid{e}_{\mathrm{2}}}}$ then there exists \ensuremath{\Varid{e'}} such that
if $\rstep{\rr}{\ensuremath{\Varid{e}_{\mathrm{1}}}}{\ensuremath{\Varid{e'}}}$ and $\rstep{\rr}{\ensuremath{\Varid{e}_{\mathrm{2}}}}{\ensuremath{\Varid{e'}}}$.
This is useful because the $\rr$-normal form (if it exists) of any expression is if $\rr$ is confluent.

But what can we say if $\rr$ is not confluent---for example, what if $\rr$ has the even-odd problem
discussed in \cref{sec:even-odd}? Ariola and Blom introduced a notion of \emph{infinite normal form},
which is a (possibly infinite) set of (possibly \emph{erased}) terms rather than a single term.
We summarize this idea, using our own terminology, as follows:
\begin{itemize}
\item Let an \emph{erasure} of a term be a copy in which one or more subterms (possibly the entire term)
  have been replaced with \ensuremath{\mathit{\Omega}}, a special term that means ``unknown'' or ``we don't know yet.''
\item We can compare possibly erased terms with the partial order $\leqomega$, which is the transitive,
  reflexive, and compatible closure of the relation in which \ensuremath{\mathit{\Omega}} is less than any other term.
  Observe that if \ensuremath{\Varid{e'}} is an erasure of \ensuremath{\Varid{e}} then $\ensuremath{\Varid{e'}} \leqomega \ensuremath{\Varid{e}}$.
\item Define the $\rr$-\emph{information content} $\omegarr(\ensuremath{\Varid{e}})$ of an expression \ensuremath{\Varid{e}} to be the unique erasure of \ensuremath{\Varid{e}} in which \emph{every} redex
  has been replaced by \ensuremath{\mathit{\Omega}}. Any non-\ensuremath{\mathit{\Omega}} structure in the result is therefore ``permanent'':
  no further reductions under $\rr$ can alter that structure.
\item For any expression \ensuremath{\Varid{e}}, consider the (possibly infinite) set $A$ of all its $\rr$-reducts, that is, \emph{all} expressions \ensuremath{\Varid{e'}} such that
  $\rstep{\rr}{e}{e'}$, then let $B = \setcomp{\omegarr(\ensuremath{\Varid{e}})}{\ensuremath{\Varid{e}} \in A}$.
  The \emph{infinite $\rr$-normal form} of \ensuremath{\Varid{e}} is a set $C$, the \emph{downward closure} of $B$ under $\leqomega$; that is,
  for every element of $B$, $C$ contains not only that element but also every possible erasure of that element.
\end{itemize}
It is the step of taking a downward closure under $\leqomega$ that makes it possible to prove that infinite normal forms
are unique for certain reduction relations.

But working with these infinite sets directly is tricky. Fortunately, there is a simpler path, because Ariola and Blom also prove
an important theorem: a reduction relation that is \emph{monotonic} (meaning that no rewrite rule ever decreases the information content)
has unique infinite normal forms if and only if the reduction relation is \emph{skew confluent} \cite[Theorem 5.4]{ariola:skew-confluence}---and skew
confluence is much easier to prove, using techniques that do not involve infinite sets, but are fairly similar to proofs of ordinary confluence
that use commutative diagrams and case analysis.

Here, then, is the idea behind skew confluence: for all \ensuremath{\Varid{e}}, \ensuremath{\Varid{e}_{\mathrm{1}}}, and \ensuremath{\Varid{e}_{\mathrm{2}}},
if $\rstep{\rr}{\ensuremath{\Varid{e}}}{\ensuremath{\Varid{e}_{\mathrm{1}}}}$ and $\rstep{\rr}{\ensuremath{\Varid{e}}}{\ensuremath{\Varid{e}_{\mathrm{2}}}}$, then there may not be any \ensuremath{\Varid{e'}} such that
if $\rstep{\rr}{\ensuremath{\Varid{e}_{\mathrm{1}}}}{\ensuremath{\Varid{e'}}}$ and $\rstep{\rr}{\ensuremath{\Varid{e}_{\mathrm{2}}}}{\ensuremath{\Varid{e'}}}$. Instead, you can choose either of \ensuremath{\Varid{e}_{\mathrm{1}}} and \ensuremath{\Varid{e}_{\mathrm{2}}} (let's say \ensuremath{\Varid{e}_{\mathrm{2}}}),
and then there does exists some \ensuremath{\Varid{e'}} such that $\rstep{\rr}{\ensuremath{\Varid{e}_{\mathrm{1}}}}{\ensuremath{\Varid{e'}}}$ and \ensuremath{\Varid{e'}} has at least as much information content as \ensuremath{\Varid{e}_{\mathrm{2}}},
that is, $\omegarr(\ensuremath{\Varid{e}_{\mathrm{2}}}) \leqomega \omegarr(\ensuremath{\Varid{e'}})$. This works both ways, by symmetry, because the choice of \ensuremath{\Varid{e}_{\mathrm{1}}} or \ensuremath{\Varid{e}_{\mathrm{2}}} is arbitrary.

In the case of the even-odd problem discussed in \cref{sec:even-odd}, further reductions on the two expressions
\ensuremath{\exists\Varid{y}.\,\Varid{y}\hspace{-0.14em}=\hspace{-0.14em}\lambda\Varid{z}.\,\langle{}\mathrm{1},\Varid{y}\rangle;\,\langle{}\mathrm{1},\Varid{y}\rangle} and \ensuremath{\exists\Varid{x}.\,\Varid{x}\mathrel{=}\langle{}\mathrm{1},\lambda\Varid{z}.\,\Varid{x}\rangle;\,\Varid{x}} might produce longer expressions such as
\ensuremath{\exists\Varid{y}.\,\Varid{y}\hspace{-0.14em}=\hspace{-0.14em}\lambda\Varid{z}.\,\langle{}\mathrm{1},\Varid{y}\rangle;\,\langle{}\mathrm{1},\lambda\Varid{z}.\,\langle{}\mathrm{1},\lambda\Varid{z}.\,\langle{}\mathrm{1},\Varid{y}\rangle\rangle\rangle}
and \ensuremath{\exists\Varid{x}.\,\Varid{x}\mathrel{=}\langle{}\mathrm{1},\lambda\Varid{z}.\,\Varid{x}\rangle;\,\langle{}\mathrm{1},\lambda\Varid{z}.\,\langle{}\mathrm{1},\lambda\Varid{z}.\,\langle{}\mathrm{1},\lambda\Varid{z}.\,\Varid{x}\rangle\rangle\rangle}.
One can see that each of these expressions
will produce a value of the form \ensuremath{\langle{}\mathrm{1},\lambda\Varid{z}.\,\langle{}\mathrm{1},\lambda\Varid{z}.\,\langle{}\mathrm{1},\cdots\rangle\rangle\rangle)}. Confluence cannot be achieved,
because no matter how many times one substitutes for \ensuremath{\Varid{y}} in the first expression or for \ensuremath{\Varid{x}} in the second expression,
the first expression will always have one more occurrence of ``\ensuremath{\mathrm{1}}'' than ``\ensuremath{\lambda\Varid{z}.\,}'' but the second expression
will always have the same number of occurrences of ``\ensuremath{\mathrm{1}}'' than ``\ensuremath{\lambda\Varid{z}.\,}''; on the other hand, no matter
how long one expression becomes through substitution, the other one can be made even longer,
so as to have at least as much information content (and in fact more, in this example). The theory of skew confluence makes this idea precise.

Another important fact is that if a reduction relation is confluent and monotonic, then it is skew confluent
\cite[Corollary 5.5]{ariola:skew-confluence}. In \cref{sec:skew-confluence} we prove a new result: if two relations are skew confluent
with respect to the same information content function and commute, then their union is also skew confluent.
(In fact, it is not required that the two relations fully commute; a slightly weaker precondition suffices.)
This provides a detailed strategy for reusing parts of the proof of confluence for \versecalc{} sketched in \cref{sec:confluence-sketch}
and detailed in \cref{sec:confluence}: our plan is to (i) define an appropriate information content function for \versecalc{} expressions;
(ii) prove that all the rewrite rules for \versecalc{} are monotonic in this information content function;
(iii) prove that the Unification rules (modified to permit recursive substitution) together with the Normalization rules are skew confluent;
(iv) prove that this combined set of rules commutes in the necessary way with the rules for Application, Elimination, and Choice (which taken together are already known to be confluent);
and (v) then apply our new result to show that the entire set of rewrite rules is skew confluent.
At this time steps (iii) and (iv) are incomplete, so we emphasize that \emph{we do not yet have a complete proof of skew confluence for \versecalc{}}.
\end{comment}

\begin{comment}

further reductions on the two terms
\ensuremath{\exists\Varid{y}.\,\Varid{y}\hspace{-0.14em}=\hspace{-0.14em}\lambda\Varid{z}.\,\langle{}\mathrm{1},\Varid{y}\rangle;\,\langle{}\mathrm{1},\Varid{y}\rangle} and \ensuremath{\exists\Varid{x}.\,\Varid{x}\mathrel{=}\langle{}\mathrm{1},\lambda\Varid{z}.\,\Varid{x}\rangle;\,\Varid{x}} might produce longer terms such as
\ensuremath{\exists\Varid{y}.\,\Varid{y}\hspace{-0.14em}=\hspace{-0.14em}\lambda\Varid{z}.\,\langle{}\mathrm{1},\Varid{y}\rangle;\,\langle{}\mathrm{1},\lambda\Varid{z}.\,\langle{}\mathrm{1},\lambda\Varid{z}.\,\langle{}\mathrm{1},\Varid{y}\rangle\rangle\rangle}
and \ensuremath{\exists\Varid{x}.\,\Varid{x}\mathrel{=}\langle{}\mathrm{1},\lambda\Varid{z}.\,\Varid{x}\rangle;\,\langle{}\mathrm{1},\lambda\Varid{z}.\,\langle{}\mathrm{1},\lambda\Varid{z}.\,\langle{}\mathrm{1},\lambda\Varid{z}.\,\Varid{x}\rangle\rangle\rangle}.
One can see that each of these terms
will produce a value of the form \ensuremath{\langle{}\mathrm{1},\lambda\Varid{z}.\,\langle{}\mathrm{1},\lambda\Varid{z}.\,\langle{}\mathrm{1},\cdots\rangle\rangle\rangle)}. Confluence cannot be achieved,
because no matter how many times one substitutes for \ensuremath{\Varid{y}} in the first expression or for \ensuremath{\Varid{x}} in the second expression,
the first expression will always have one more occurrence of ``\ensuremath{\mathrm{1}}'' than ``\ensuremath{\lambda\Varid{z}.\,}'' but the second expression
will always have the same number of occurrences of ``\ensuremath{\mathrm{1}}'' than ``\ensuremath{\lambda\Varid{z}.\,}''; on the other hand, no matter
how long one expression becomes through substitution, the other one can be made even longer,
so as to have at least as much information content (and in fact more, in this example). The theory of skew confluence makes this idea precise.

Old...
and either term can be further reduced so as to contain \emph{all} the permanent structure of the other.
This provides a notion of convergence even for infinite reduction sequences.

In the case of the even-odd problem discussed in \cref{sec:even-odd}, further reductions on the two terms
\ensuremath{\exists\Varid{y}.\,\Varid{y}\hspace{-0.14em}=\hspace{-0.14em}\lambda\Varid{z}.\,\langle{}\mathrm{1},\Varid{y}\rangle;\,\langle{}\mathrm{1},\Varid{y}\rangle} and \ensuremath{\exists\Varid{x}.\,\Varid{x}\mathrel{=}\langle{}\mathrm{1},\lambda\Varid{z}.\,\Varid{x}\rangle;\,\Varid{x}} might produce longer terms such as
\ensuremath{\exists\Varid{y}.\,\Varid{y}\hspace{-0.14em}=\hspace{-0.14em}\lambda\Varid{z}.\,\langle{}\mathrm{1},\Varid{y}\rangle;\,\langle{}\mathrm{1},\lambda\Varid{z}.\,\langle{}\mathrm{1},\lambda\Varid{z}.\,\langle{}\mathrm{1},\Varid{y}\rangle\rangle\rangle}
and \ensuremath{\exists\Varid{x}.\,\Varid{x}\mathrel{=}\langle{}\mathrm{1},\lambda\Varid{z}.\,\Varid{x}\rangle;\,\langle{}\mathrm{1},\lambda\Varid{z}.\,\langle{}\mathrm{1},\lambda\Varid{z}.\,\langle{}\mathrm{1},\lambda\Varid{z}.\,\Varid{x}\rangle\rangle\rangle}.
One can see that each of these terms
will produce a value of the form \ensuremath{\langle{}\mathrm{1},\lambda\Varid{z}.\,\langle{}\mathrm{1},\lambda\Varid{z}.\,\langle{}\mathrm{1},\cdots\rangle\rangle\rangle)}. Confluence cannot be achieved,
because no matter how many times one substitutes for \ensuremath{\Varid{y}} in the first expression or for \ensuremath{\Varid{x}} in the second expression,
the first expression will always have one more occurrence of ``\ensuremath{\mathrm{1}}'' than ``\ensuremath{\lambda\Varid{z}.\,}'' but the second expression
will always have the same number of occurrences of ``\ensuremath{\mathrm{1}}'' than ``\ensuremath{\lambda\Varid{z}.\,}''; on the other hand, no matter
how long one expression becomes through substitution, the other one can be made even longer,
so as to have at least as much information content (and in fact more, in this example). The theory of skew confluence makes this idea precise.
\end{comment}


Why use skew confluence? Ordinary confluence is useful because if term \ensuremath{\Varid{e}} has an $\rr$-normal form, then that normal form is \emph{unique} if $\rr$ is confluent.
Ariola and Blom define a related notion, which we will refer to as \emph{$\rr$-skew-normal form}\footnote{Ariola and Blom call it the ``infinite normal form''; this is a bit misleading because in fact not all such forms are infinite.}, and prove that under certain conditions, if term \ensuremath{\Varid{e}} has an $\rr$-skew-normal form, then that normal form is unique if $\rr$ is skew confluent.

A  $\rr$-skew-normal form is not a single term, but rather a possibly infinite set of \emph{erased} terms.
We summarize this idea, using our own terminology, as follows:
\begin{itemize}
\item Let an \emph{erasure} of a term be a copy in which some number of subterms (possibly zero, and possibly the entire term)
  have been replaced with \ensuremath{\mathit{\Omega}}, a special term that means ``unknown'' or ``we don't know yet.''
\item We can compare erased terms with the partial order $\leqomega$, which is the transitive,
  reflexive, and compatible closure of the relation in which \ensuremath{\mathit{\Omega}} is less than any other term.
  Observe that if \ensuremath{\Varid{e'}} is any erasure of \ensuremath{\Varid{e}} (including \ensuremath{\Varid{e}} itself) then $\ensuremath{\Varid{e'}} \leqomega \ensuremath{\Varid{e}}$.
\item Define the $\rr$-\emph{information content} $\omegarr(\ensuremath{\Varid{e}})$ of a term \ensuremath{\Varid{e}} to be the unique erasure of \ensuremath{\Varid{e}} in which \emph{every} redex
  has been replaced by \ensuremath{\mathit{\Omega}}. Any non-\ensuremath{\mathit{\Omega}} structure in the result is therefore ``permanent'':
  no further reductions under $\rr$ can alter that structure.
\item Define the \emph{downward closure} $\downclose A$ of a set of terms $A$ to be the set of all elements of $A$ and all possible erasures of those elements,
  that is, $\downclose A = \setcomp{\ensuremath{\Varid{e'}}}{\ensuremath{\Varid{e}} \in A, \ensuremath{\Varid{e'}} \leqomega}$.
\item
  Define the \emph{$\rr$-skew-normal form} of \ensuremath{\Varid{e}} to be the downward closure of the set of information contents of all possible $\rr$-reducts of \ensuremath{\Varid{e}},
  that is, $\downclose \setcomp{\omega(\ensuremath{\Varid{e'}})}{\rsteps{\rr}{e}{e'}}$.
\end{itemize}
Taking the downward closure with respect to $\leqomega$ is crucial; without that step, it would not be possible to prove that skew-normal forms
are unique for certain reduction relations.

Skew confluence is a natural extension of confluence, in this sense: if $\rr$ is skew confluent,
then an expression \ensuremath{\Varid{e}} has a unique $\rr$-normal form
if and only if its $\rr$-skew-normal form is the (finite) set consisting of all possible erasures of that $\rr$-normal form.
(For example, \ensuremath{\langle{}\mathrm{1},\mathrm{2}\rangle} is the unique normal form of \ensuremath{\exists\Varid{x}.\,\Varid{x}\mathrel{=}\mathrm{2};\,\langle{}\mathrm{1},\Varid{x}\rangle}, and the  $\rr$-skew-normal form of that same term
is $\{ \ensuremath{\langle{}\mathrm{1},\mathrm{2}\rangle}, \ensuremath{\langle{}\mathit{\Omega},\mathrm{2}\rangle}, \ensuremath{\langle{}\mathrm{1},\mathit{\Omega}\rangle}, \ensuremath{\langle{}\mathit{\Omega},\mathit{\Omega}\rangle}, \ensuremath{\mathit{\Omega}} \}$.)

But working with possibly infinite sets directly is tricky. Fortunately, there is a simpler path, because Ariola and Blom prove
an important theorem:
Define the partial order $\ensuremath{\Varid{e}} \preceqomegarr \ensuremath{\Varid{e'}}$ to mean $\omegarr(\ensuremath{\Varid{e}})\leqomega\omegarr(\ensuremath{\Varid{e'}})$; then
a reduction relation that is \emph{monotonic} in $\preceqomegarr$ ($\rstep{\rr}{e}{e'} \implies e \preceqomegarr e'$)
has unique skew-normal forms if and only if the reduction relation is skew confluent \cite[Theorem 5.4]{ariola:skew-confluence}---and skew
confluence is much easier to prove, using techniques that do not involve possibly infinite sets, but are fairly similar to proofs of ordinary confluence
that use commutative diagrams and case analysis. They also prove that if a reduction relation is confluent and monotonic, then it is skew confluent
\cite[Corollary 5.5]{ariola:skew-confluence}.




\begin{definition}
  Reduction relation $\rr$ over the set of terms $T$ is \emph{skew confluent} using quasi order $\preceq$ if for all $a, b, c \in T$,
  if $a \gos{\rr} b$ and $a \gos{\rr} c$, there exists $d\in T$ such that $b \gos{\rr} d$ and $c \preceq d$. As a diagram:
  
% https://q.uiver.app/?q=WzAsNCxbMCwwLCJhIl0sWzEsMCwiYiJdLFsxLDEsImQiXSxbMCwxLCJjIl0sWzAsMSwiXFxyciIsMCx7InN0eWxlIjp7ImhlYWQiOnsibmFtZSI6ImVwaSJ9fX1dLFsxLDIsIlxccnIiLDAseyJzdHlsZSI6eyJib2R5Ijp7Im5hbWUiOiJkb3R0ZWQifSwiaGVhZCI6eyJuYW1lIjoiZXBpIn19fV0sWzAsMywiXFxyciIsMix7InN0eWxlIjp7ImhlYWQiOnsibmFtZSI6ImVwaSJ9fX1dLFszLDIsIlxccnIiLDIseyJzdHlsZSI6eyJib2R5Ijp7Im5hbWUiOiJkb3R0ZWQifSwiaGVhZCI6eyJuYW1lIjoiZXBpIn19fV0sWzIsMywiXFxwcmVjZXEiLDIseyJzdHlsZSI6eyJib2R5Ijp7Im5hbWUiOiJub25lIn0sImhlYWQiOnsibmFtZSI6ImVwaSJ9fX1dLFs2LDUsIiIsMSx7InNob3J0ZW4iOnsic291cmNlIjoyMCwidGFyZ2V0IjoyMH0sInN0eWxlIjp7ImJvZHkiOnsibmFtZSI6Im5vbmUifSwiaGVhZCI6eyJuYW1lIjoibm9uZSJ9fX1dXQ==
\[\begin{tikzcd}
	a & b \\
	c & d
	\arrow["\rr", two heads, from=1-1, to=1-2]
	\arrow[""{name=0, anchor=center, inner sep=0}, "\rr", dotted, two heads, from=1-2, to=2-2]
	\arrow[""{name=1, anchor=center, inner sep=0}, "\rr"', two heads, from=1-1, to=2-1]
	\arrow["\rr"', dotted, two heads, from=2-1, to=2-2]
	\arrow["\preceq"', no body, two heads, from=2-2, to=2-1]
	\arrow[draw=none, from=1, to=0]
\end{tikzcd}\]
\end{definition}

\skewtodo{More to come.}

\subsubsection{New Lemmas about Skew Confluence}

\begin{lemma}
If relation $\rr$ is monotonic in some quasi order $\preceq$ and confluent, then it is skew confluent using $\preceq$.
\end{lemma}
\begin{proof}
  By the definition of confluence,
  % https://q.uiver.app/?q=WzAsNCxbMCwwLCJhIl0sWzEsMCwiYyJdLFswLDEsImIiXSxbMSwxLCJkIl0sWzAsMiwiXFxyciIsMix7InN0eWxlIjp7ImhlYWQiOnsibmFtZSI6ImVwaSJ9fX1dLFswLDEsIlxccnIiLDAseyJzdHlsZSI6eyJoZWFkIjp7Im5hbWUiOiJlcGkifX19XSxbMSwzLCJcXHJyIiwwLHsic3R5bGUiOnsiYm9keSI6eyJuYW1lIjoiZG90dGVkIn0sImhlYWQiOnsibmFtZSI6ImVwaSJ9fX1dLFsyLDMsIlxccnIiLDIseyJzdHlsZSI6eyJib2R5Ijp7Im5hbWUiOiJkb3R0ZWQifSwiaGVhZCI6eyJuYW1lIjoiZXBpIn19fV1d
\[\begin{tikzcd}
	a & c \\
	b & d
	\arrow["\rr"', two heads, from=1-1, to=2-1]
	\arrow["\rr", two heads, from=1-1, to=1-2]
	\arrow["\rr", dotted, two heads, from=1-2, to=2-2]
	\arrow["\rr"', dotted, two heads, from=2-1, to=2-2]
\end{tikzcd}\]
Because $\rr$ is monotonic, $\gos{\rr} \subset \inccon{\preceq}{\rr}$, therefore 
% https://q.uiver.app/?q=WzAsNCxbMCwwLCJhIl0sWzEsMCwiYyJdLFswLDEsImIiXSxbMSwxLCJkIl0sWzAsMiwiXFxyciIsMix7InN0eWxlIjp7ImhlYWQiOnsibmFtZSI6ImVwaSJ9fX1dLFswLDEsIlxccnIiLDAseyJzdHlsZSI6eyJoZWFkIjp7Im5hbWUiOiJlcGkifX19XSxbMSwzLCJcXHJyIiwwLHsic3R5bGUiOnsiYm9keSI6eyJuYW1lIjoiZG90dGVkIn0sImhlYWQiOnsibmFtZSI6ImVwaSJ9fX1dLFsyLDMsIlxccnIiLDIseyJzdHlsZSI6eyJib2R5Ijp7Im5hbWUiOiJkb3R0ZWQifSwiaGVhZCI6eyJuYW1lIjoiZXBpIn19fV0sWzMsMiwiXFxwcmVjZXEiLDIseyJzdHlsZSI6eyJib2R5Ijp7Im5hbWUiOiJub25lIn0sImhlYWQiOnsibmFtZSI6ImVwaSJ9fX1dXQ==
\[\begin{tikzcd}
	a & c \\
	b & d
	\arrow["\rr"', two heads, from=1-1, to=2-1]
	\arrow["\rr", two heads, from=1-1, to=1-2]
	\arrow["\rr", dotted, two heads, from=1-2, to=2-2]
	\arrow["\rr"', dotted, two heads, from=2-1, to=2-2]
	\arrow["\preceq"', no body, two heads, from=2-2, to=2-1]
\end{tikzcd}\]
\end{proof}

\begin{definition}
  Let reduction relation $\rr$ be monotonic in quasi order $\preceq$ and skew confluent using $\preceq$.
  Let reduction relation $\rr_\preceq^{\leftarrow}$ be defined by a set of rewrite rules that are converses of those rewrite rules of $\rr$ whose converses are used in the proof that $\rr$ is skew
  confluent.
  Then we say that $\rr$ is \emph{skew confluent using $\preceq$ and  $\rr_\preceq^{\leftarrow}$}.
\end{definition}


\begin{lemma}\label{lem:four-diagrams}
  Let relation $\rr$ be monotonic in quasi order $\preceq$ and skew confluent using $\preceq$ and $\rr_\preceq^{\leftarrow}$;
  similarly let $\rs$ be monotonic in the same quasi order $\preceq$ and skew confluent using $\preceq$ and $\rs_\preceq^{\leftarrow}$.
  Suppose furthermore that $\rr$ and $\rs$ commute and that $\rr$ and $\rs_\preceq^{\leftarrow}$ commute. 
  Then the following four diagrams hold:

% https://q.uiver.app/?q=WzAsNCxbMCwwLCJhIl0sWzEsMCwiYiJdLFsxLDEsImQiXSxbMCwxLCJjIl0sWzAsMSwiXFxyciIsMl0sWzEsMiwiXFxyciIsMCx7InN0eWxlIjp7ImJvZHkiOnsibmFtZSI6ImRvdHRlZCJ9LCJoZWFkIjp7Im5hbWUiOiJlcGkifX19XSxbMCwzLCJcXHJyIiwyLHsic3R5bGUiOnsiaGVhZCI6eyJuYW1lIjoiZXBpIn19fV0sWzMsMiwiXFxyclxcY3VwXFxycyIsMix7InN0eWxlIjp7ImJvZHkiOnsibmFtZSI6ImRvdHRlZCJ9LCJoZWFkIjp7Im5hbWUiOiJlcGkifX19XSxbMiwzLCJcXHByZWNlcSIsMix7InN0eWxlIjp7ImJvZHkiOnsibmFtZSI6Im5vbmUifSwiaGVhZCI6eyJuYW1lIjoiZXBpIn19fV0sWzYsNSwiIiwxLHsic2hvcnRlbiI6eyJzb3VyY2UiOjIwLCJ0YXJnZXQiOjIwfSwic3R5bGUiOnsiYm9keSI6eyJuYW1lIjoibm9uZSJ9LCJoZWFkIjp7Im5hbWUiOiJub25lIn19fV1d
\[\begin{tikzcd}
	a & b \\
	c & d
	\arrow["\rr"', from=1-1, to=1-2]
	\arrow[""{name=0, anchor=center, inner sep=0}, "\rr", dotted, two heads, from=1-2, to=2-2]
	\arrow[""{name=1, anchor=center, inner sep=0}, "\rr"', two heads, from=1-1, to=2-1]
	\arrow["\rr\cup\rs"', dotted, two heads, from=2-1, to=2-2]
	\arrow["\preceq"', no body, two heads, from=2-2, to=2-1]
	\arrow[draw=none, from=1, to=0]
      \end{tikzcd}
      \qquad\qquad
% https://q.uiver.app/?q=WzAsNCxbMCwwLCJhIl0sWzEsMCwiYiJdLFsxLDEsImQiXSxbMCwxLCJjIl0sWzAsMSwiXFxycyIsMl0sWzEsMiwiXFxyciIsMCx7InN0eWxlIjp7ImJvZHkiOnsibmFtZSI6ImRvdHRlZCJ9LCJoZWFkIjp7Im5hbWUiOiJlcGkifX19XSxbMCwzLCJcXHJyIiwyLHsic3R5bGUiOnsiaGVhZCI6eyJuYW1lIjoiZXBpIn19fV0sWzMsMiwiXFxyclxcY3VwXFxycyIsMix7InN0eWxlIjp7ImJvZHkiOnsibmFtZSI6ImRvdHRlZCJ9LCJoZWFkIjp7Im5hbWUiOiJlcGkifX19XSxbMiwzLCJcXHByZWNlcSIsMix7InN0eWxlIjp7ImJvZHkiOnsibmFtZSI6Im5vbmUifSwiaGVhZCI6eyJuYW1lIjoiZXBpIn19fV0sWzYsNSwiIiwxLHsic2hvcnRlbiI6eyJzb3VyY2UiOjIwLCJ0YXJnZXQiOjIwfSwic3R5bGUiOnsiYm9keSI6eyJuYW1lIjoibm9uZSJ9LCJoZWFkIjp7Im5hbWUiOiJub25lIn19fV1d
\begin{tikzcd}
	a & b \\
	c & d
	\arrow["\rs"', from=1-1, to=1-2]
	\arrow[""{name=0, anchor=center, inner sep=0}, "\rr", dotted, two heads, from=1-2, to=2-2]
	\arrow[""{name=1, anchor=center, inner sep=0}, "\rr"', two heads, from=1-1, to=2-1]
	\arrow["\rr\cup\rs"', dotted, two heads, from=2-1, to=2-2]
	\arrow["\preceq"', no body, two heads, from=2-2, to=2-1]
	\arrow[draw=none, from=1, to=0]
      \end{tikzcd}
      \qquad\qquad
    % https://q.uiver.app/?q=WzAsNCxbMCwwLCJhIl0sWzEsMCwiYiJdLFsxLDEsImQiXSxbMCwxLCJjIl0sWzEsMCwiXFxyciJdLFsxLDIsIlxccnIiLDAseyJzdHlsZSI6eyJib2R5Ijp7Im5hbWUiOiJkb3R0ZWQifSwiaGVhZCI6eyJuYW1lIjoiZXBpIn19fV0sWzAsMywiXFxyciIsMix7InN0eWxlIjp7ImhlYWQiOnsibmFtZSI6ImVwaSJ9fX1dLFszLDIsIlxccnJcXGN1cFxccnMiLDIseyJzdHlsZSI6eyJib2R5Ijp7Im5hbWUiOiJkb3R0ZWQifSwiaGVhZCI6eyJuYW1lIjoiZXBpIn19fV0sWzIsMywiXFxwcmVjZXEiLDIseyJzdHlsZSI6eyJib2R5Ijp7Im5hbWUiOiJub25lIn0sImhlYWQiOnsibmFtZSI6ImVwaSJ9fX1dLFs2LDUsIiIsMSx7InNob3J0ZW4iOnsic291cmNlIjoyMCwidGFyZ2V0IjoyMH0sInN0eWxlIjp7ImJvZHkiOnsibmFtZSI6Im5vbmUifSwiaGVhZCI6eyJuYW1lIjoibm9uZSJ9fX1dXQ==
\begin{tikzcd}
	a & b \\
	c & d
	\arrow["\rr", from=1-2, to=1-1]
	\arrow[""{name=0, anchor=center, inner sep=0}, "\rr", dotted, two heads, from=1-2, to=2-2]
	\arrow[""{name=1, anchor=center, inner sep=0}, "\rr"', two heads, from=1-1, to=2-1]
	\arrow["\rr\cup\rs"', dotted, two heads, from=2-1, to=2-2]
	\arrow["\preceq"', no body, two heads, from=2-2, to=2-1]
	\arrow[draw=none, from=1, to=0]
      \end{tikzcd}
      \qquad\qquad
% https://q.uiver.app/?q=WzAsNCxbMCwwLCJhIl0sWzEsMCwiYiJdLFsxLDEsImQiXSxbMCwxLCJjIl0sWzEsMCwiXFxycyJdLFsxLDIsIlxccnIiLDAseyJzdHlsZSI6eyJib2R5Ijp7Im5hbWUiOiJkb3R0ZWQifSwiaGVhZCI6eyJuYW1lIjoiZXBpIn19fV0sWzAsMywiXFxyciIsMix7InN0eWxlIjp7ImhlYWQiOnsibmFtZSI6ImVwaSJ9fX1dLFszLDIsIlxccnJcXGN1cFxccnMiLDIseyJzdHlsZSI6eyJib2R5Ijp7Im5hbWUiOiJkb3R0ZWQifSwiaGVhZCI6eyJuYW1lIjoiZXBpIn19fV0sWzIsMywiXFxwcmVjZXEiLDIseyJzdHlsZSI6eyJib2R5Ijp7Im5hbWUiOiJub25lIn0sImhlYWQiOnsibmFtZSI6ImVwaSJ9fX1dLFs2LDUsIiIsMSx7InNob3J0ZW4iOnsic291cmNlIjoyMCwidGFyZ2V0IjoyMH0sInN0eWxlIjp7ImJvZHkiOnsibmFtZSI6Im5vbmUifSwiaGVhZCI6eyJuYW1lIjoibm9uZSJ9fX1dXQ==
\begin{tikzcd}
	a & b \\
	c & d
	\arrow["\rs", from=1-2, to=1-1]
	\arrow[""{name=0, anchor=center, inner sep=0}, "\rr", dotted, two heads, from=1-2, to=2-2]
	\arrow[""{name=1, anchor=center, inner sep=0}, "\rr"', two heads, from=1-1, to=2-1]
	\arrow["\rr\cup\rs"', dotted, two heads, from=2-1, to=2-2]
	\arrow["\preceq"', no body, two heads, from=2-2, to=2-1]
	\arrow[draw=none, from=1, to=0]
      \end{tikzcd}\]
  \end{lemma}
  \begin{proof}
    
    (1) Because $\rr$ is skew confluent, we have

    % https://q.uiver.app/?q=WzAsNCxbMCwwLCJhIl0sWzEsMCwiYiJdLFsxLDEsImQiXSxbMCwxLCJjIl0sWzAsMSwiXFxyciIsMl0sWzEsMiwiXFxyciIsMCx7InN0eWxlIjp7ImJvZHkiOnsibmFtZSI6ImRvdHRlZCJ9LCJoZWFkIjp7Im5hbWUiOiJlcGkifX19XSxbMCwzLCJcXHJyIiwyLHsic3R5bGUiOnsiaGVhZCI6eyJuYW1lIjoiZXBpIn19fV0sWzMsMiwiXFxyciIsMix7InN0eWxlIjp7ImJvZHkiOnsibmFtZSI6ImRvdHRlZCJ9LCJoZWFkIjp7Im5hbWUiOiJlcGkifX19XSxbMiwzLCJcXHByZWNlcSIsMix7InN0eWxlIjp7ImJvZHkiOnsibmFtZSI6Im5vbmUifSwiaGVhZCI6eyJuYW1lIjoiZXBpIn19fV0sWzYsNSwiIiwxLHsic2hvcnRlbiI6eyJzb3VyY2UiOjIwLCJ0YXJnZXQiOjIwfSwic3R5bGUiOnsiYm9keSI6eyJuYW1lIjoibm9uZSJ9LCJoZWFkIjp7Im5hbWUiOiJub25lIn19fV1d
\[\begin{tikzcd}
	a & b \\
	c & d
	\arrow["\rr"', from=1-1, to=1-2]
	\arrow[""{name=0, anchor=center, inner sep=0}, "\rr", dotted, two heads, from=1-2, to=2-2]
	\arrow[""{name=1, anchor=center, inner sep=0}, "\rr"', two heads, from=1-1, to=2-1]
	\arrow["\rr"', dotted, two heads, from=2-1, to=2-2]
	\arrow["\preceq"', no body, two heads, from=2-2, to=2-1]
	\arrow[draw=none, from=1, to=0]
      \end{tikzcd}\]

\noindent    Because $\inccon{\preceq}{\rr} \subset \inccon{\preceq}{\rr\cup\rs}$, the first diagram follows.

        (2) Because $\rr$ and $\rs$ commute, we have

% https://q.uiver.app/?q=WzAsNCxbMCwwLCJhIl0sWzEsMCwiYiJdLFsxLDEsImQiXSxbMCwxLCJjIl0sWzAsMSwiXFxycyIsMl0sWzEsMiwiXFxyciIsMCx7InN0eWxlIjp7ImJvZHkiOnsibmFtZSI6ImRvdHRlZCJ9LCJoZWFkIjp7Im5hbWUiOiJlcGkifX19XSxbMCwzLCJcXHJyIiwyLHsic3R5bGUiOnsiaGVhZCI6eyJuYW1lIjoiZXBpIn19fV0sWzMsMiwiXFxycyIsMix7InN0eWxlIjp7ImJvZHkiOnsibmFtZSI6ImRvdHRlZCJ9LCJoZWFkIjp7Im5hbWUiOiJlcGkifX19XSxbNiw1LCIiLDEseyJzaG9ydGVuIjp7InNvdXJjZSI6MjAsInRhcmdldCI6MjB9LCJzdHlsZSI6eyJib2R5Ijp7Im5hbWUiOiJub25lIn0sImhlYWQiOnsibmFtZSI6Im5vbmUifX19XV0=
\[\begin{tikzcd}
	a & b \\
	c & d
	\arrow["\rs"', from=1-1, to=1-2]
	\arrow[""{name=0, anchor=center, inner sep=0}, "\rr", dotted, two heads, from=1-2, to=2-2]
	\arrow[""{name=1, anchor=center, inner sep=0}, "\rr"', two heads, from=1-1, to=2-1]
	\arrow["\rs"', dotted, two heads, from=2-1, to=2-2]
	\arrow[draw=none, from=1, to=0]
\end{tikzcd}\]

\noindent    Because $\gos{\rs} \subset \inccon{\preceq}{\rr\cup\rs}$, the second diagram follows.

(3) The following diagram clearly holds if the sequence of reduction steps from \ensuremath{\Varid{b}} to \ensuremath{\Varid{d}} is the same as the sequence of reduction steps from \ensuremath{\Varid{b}} to \ensuremath{\Varid{a}} to \ensuremath{\Varid{d}}:

% https://q.uiver.app/?q=WzAsNCxbMCwwLCJhIl0sWzEsMCwiYiJdLFsxLDEsImQiXSxbMCwxLCJjIl0sWzEsMCwiXFxyciJdLFsxLDIsIlxccnIiLDAseyJzdHlsZSI6eyJib2R5Ijp7Im5hbWUiOiJkb3R0ZWQifSwiaGVhZCI6eyJuYW1lIjoiZXBpIn19fV0sWzAsMywiXFxyciIsMix7InN0eWxlIjp7ImhlYWQiOnsibmFtZSI6ImVwaSJ9fX1dLFszLDIsIlxcZXF1aXYiLDEseyJzdHlsZSI6eyJib2R5Ijp7Im5hbWUiOiJub25lIn0sImhlYWQiOnsibmFtZSI6Im5vbmUifX19XSxbNiw1LCIiLDEseyJzaG9ydGVuIjp7InNvdXJjZSI6MjAsInRhcmdldCI6MjB9LCJzdHlsZSI6eyJib2R5Ijp7Im5hbWUiOiJub25lIn0sImhlYWQiOnsibmFtZSI6Im5vbmUifX19XV0=
\[\begin{tikzcd}
	a & b \\
	c & d
	\arrow["\rr", from=1-2, to=1-1]
	\arrow[""{name=0, anchor=center, inner sep=0}, "\rr", dotted, two heads, from=1-2, to=2-2]
	\arrow[""{name=1, anchor=center, inner sep=0}, "\rr"', two heads, from=1-1, to=2-1]
	\arrow["\equiv"{description}, draw=none, from=2-1, to=2-2]
	\arrow[draw=none, from=1, to=0]
      \end{tikzcd}\]

    \noindent    Because $\equiv \subset \inccon{\preceq}{\rr\cup\rs}$, the third diagram follows.

    (4) Because $\rr$ and $\rs_\preceq^{\leftarrow}$ commute, we have

    % https://q.uiver.app/?q=WzAsNCxbMCwwLCJhIl0sWzEsMCwiYiJdLFsxLDEsImQiXSxbMCwxLCJjIl0sWzEsMCwiXFxycyJdLFsxLDIsIlxccnIiLDAseyJzdHlsZSI6eyJib2R5Ijp7Im5hbWUiOiJkb3R0ZWQifSwiaGVhZCI6eyJuYW1lIjoiZXBpIn19fV0sWzAsMywiXFxyciIsMix7InN0eWxlIjp7ImhlYWQiOnsibmFtZSI6ImVwaSJ9fX1dLFsyLDMsIlxccnMiLDAseyJzdHlsZSI6eyJib2R5Ijp7Im5hbWUiOiJkb3R0ZWQifSwiaGVhZCI6eyJuYW1lIjoiZXBpIn19fV0sWzYsNSwiIiwxLHsic2hvcnRlbiI6eyJzb3VyY2UiOjIwLCJ0YXJnZXQiOjIwfSwic3R5bGUiOnsiYm9keSI6eyJuYW1lIjoibm9uZSJ9LCJoZWFkIjp7Im5hbWUiOiJub25lIn19fV1d
\[\begin{tikzcd}
	a & b \\
	c & d
	\arrow["\rs", from=1-2, to=1-1]
	\arrow[""{name=0, anchor=center, inner sep=0}, "\rr", dotted, two heads, from=1-2, to=2-2]
	\arrow[""{name=1, anchor=center, inner sep=0}, "\rr"', two heads, from=1-1, to=2-1]
	\arrow["\rs", dotted, two heads, from=2-2, to=2-1]
	\arrow[draw=none, from=1, to=0]
      \end{tikzcd}\]

        \noindent    Because $\ungos{\rs} \subset \inccon{\preceq}{\rr\cup\rs}$, the fourth diagram follows.

        \skewtodo{It may turn out to be impossible to prove for our specific application that $\rr$ and $\rs_\preceq^{\leftarrow}$ commute.
          In that case, it may be necessary to use a more complicated or more subtle precondition.
        The important thing is to prove the fourth diagram somehow.}
\end{proof}






  
\begin{lemma}\label{lem:skewrowinduction}
  Let relation $\rr$ be monotonic in quasi order $\preceq$ and skew confluent using $\preceq$ and $\rr_\preceq^{\leftarrow}$;
  similarly let $\rs$ be monotonic in the same quasi order $\preceq$ and skew confluent using $\preceq$ and $\rs_\preceq^{\leftarrow}$.
  Suppose furthermore that $\rr$ and $\rs$ commute and that $\rr$ and $\rs_\preceq^{\leftarrow}$ commute. 
  Then

  % https://q.uiver.app/?q=WzAsNCxbMCwwLCJhIl0sWzEsMCwiYiJdLFsxLDEsImQiXSxbMCwxLCJjIl0sWzAsMSwiXFxyclxcY3VwXFxycyIsMix7InN0eWxlIjp7ImJvZHkiOnsibmFtZSI6ImRvdHRlZCJ9LCJoZWFkIjp7Im5hbWUiOiJlcGkifX19XSxbMSwwLCJcXHByZWNlcSIsMix7InN0eWxlIjp7ImJvZHkiOnsibmFtZSI6Im5vbmUifSwiaGVhZCI6eyJuYW1lIjoiZXBpIn19fV0sWzEsMiwiXFxyciIsMCx7InN0eWxlIjp7ImJvZHkiOnsibmFtZSI6ImRvdHRlZCJ9LCJoZWFkIjp7Im5hbWUiOiJlcGkifX19XSxbMCwzLCJcXHJyIiwyLHsic3R5bGUiOnsiaGVhZCI6eyJuYW1lIjoiZXBpIn19fV0sWzMsMiwiXFxyclxcY3VwXFxycyIsMix7InN0eWxlIjp7ImJvZHkiOnsibmFtZSI6ImRvdHRlZCJ9LCJoZWFkIjp7Im5hbWUiOiJlcGkifX19XSxbMiwzLCJcXHByZWNlcSIsMix7InN0eWxlIjp7ImJvZHkiOnsibmFtZSI6Im5vbmUifSwiaGVhZCI6eyJuYW1lIjoiZXBpIn19fV0sWzcsNiwiIiwxLHsic2hvcnRlbiI6eyJzb3VyY2UiOjIwLCJ0YXJnZXQiOjIwfSwic3R5bGUiOnsiYm9keSI6eyJuYW1lIjoibm9uZSJ9LCJoZWFkIjp7Im5hbWUiOiJub25lIn19fV1d
\[\begin{tikzcd}
	a & b \\
	c & d
	\arrow["\rr\cup\rs"', dotted, two heads, from=1-1, to=1-2]
	\arrow["\preceq"', no body, two heads, from=1-2, to=1-1]
	\arrow[""{name=0, anchor=center, inner sep=0}, "\rr", dotted, two heads, from=1-2, to=2-2]
	\arrow[""{name=1, anchor=center, inner sep=0}, "\rr"', two heads, from=1-1, to=2-1]
	\arrow["\rr\cup\rs"', dotted, two heads, from=2-1, to=2-2]
	\arrow["\preceq"', no body, two heads, from=2-2, to=2-1]
	\arrow[draw=none, from=1, to=0]
\end{tikzcd}\]
\end{lemma}
\begin{proof}
  By induction on the size of the top edge of the diagram. At each step one of the four diagrams from \cref{lem:four-diagrams} will be used.

  
 \skewtodo{More to come.}
\end{proof}


\begin{lemma}\label{lem:skewrow}
  Let relation $\rr$ be monotonic in quasi order $\preceq$ and skew confluent using $\preceq$ and $\rr_\preceq^\leftarrow$;
  similarly let $\rs$ be monotonic in the same quasi order $\preceq$ and skew confluent using $\preceq$ and $\rs_\preceq^\leftarrow$.
  Then

  % https://q.uiver.app/?q=WzAsNCxbMCwwLCJhIl0sWzEsMCwiYiJdLFsxLDEsImQiXSxbMCwxLCJjIl0sWzAsMSwiXFxyclxcY3VwXFxycyIsMix7InN0eWxlIjp7ImJvZHkiOnsibmFtZSI6ImRvdHRlZCJ9LCJoZWFkIjp7Im5hbWUiOiJlcGkifX19XSxbMSwwLCJcXHByZWNlcSIsMix7InN0eWxlIjp7ImJvZHkiOnsibmFtZSI6Im5vbmUifSwiaGVhZCI6eyJuYW1lIjoiZXBpIn19fV0sWzEsMiwiXFxyclxcY3VwXFxycyIsMCx7InN0eWxlIjp7ImJvZHkiOnsibmFtZSI6ImRvdHRlZCJ9LCJoZWFkIjp7Im5hbWUiOiJlcGkifX19XSxbMCwzLCJcXHJyXFxjdXBcXHJzIiwyXSxbMywyLCJcXHJyXFxjdXBcXHJzIiwyLHsic3R5bGUiOnsiYm9keSI6eyJuYW1lIjoiZG90dGVkIn0sImhlYWQiOnsibmFtZSI6ImVwaSJ9fX1dLFsyLDMsIlxccHJlY2VxIiwyLHsic3R5bGUiOnsiYm9keSI6eyJuYW1lIjoibm9uZSJ9LCJoZWFkIjp7Im5hbWUiOiJlcGkifX19XSxbNyw2LCIiLDEseyJzaG9ydGVuIjp7InNvdXJjZSI6MjAsInRhcmdldCI6MjB9LCJzdHlsZSI6eyJib2R5Ijp7Im5hbWUiOiJub25lIn0sImhlYWQiOnsibmFtZSI6Im5vbmUifX19XV0=
\[\begin{tikzcd}
	a & b \\
	c & d
	\arrow["\rr\cup\rs"', dotted, two heads, from=1-1, to=1-2]
	\arrow["\preceq"', no body, two heads, from=1-2, to=1-1]
	\arrow[""{name=0, anchor=center, inner sep=0}, "\rr\cup\rs", dotted, two heads, from=1-2, to=2-2]
	\arrow[""{name=1, anchor=center, inner sep=0}, "\rr\cup\rs"', from=1-1, to=2-1]
	\arrow["\rr\cup\rs"', dotted, two heads, from=2-1, to=2-2]
	\arrow["\preceq"', no body, two heads, from=2-2, to=2-1]
	\arrow[draw=none, from=1, to=0]
\end{tikzcd}\]

\end{lemma}
\begin{proof}
By case analysis on whether left edge uses $\rr$ or $\rs$; then project that left edge into $\rr^*$ or $\rs^*$ respectively and apply \cref{lem:skewrowinduction}.

  
 \skewtodo{More to come.}
\end{proof}


\begin{lemma}\label{lem:skewgrid}
  Let relation $\rr$ be monotonic in quasi order $\preceq$ and skew confluent using $\preceq$ and $\rr_\preceq^\leftarrow$;
  similarly let $\rs$ be monotonic in the same quasi order $\preceq$ and skew confluent using $\preceq$ and $\rs_\preceq^\leftarrow$.
  Then

% https://q.uiver.app/?q=WzAsNCxbMCwwLCJhIl0sWzEsMCwiYiJdLFswLDEsImMiXSxbMSwxLCJkIl0sWzAsMiwiXFxyclxcY3VwXFxycyIsMix7InN0eWxlIjp7ImhlYWQiOnsibmFtZSI6ImVwaSJ9fX1dLFswLDEsIlxccnJcXGN1cFxccnMiLDIseyJzdHlsZSI6eyJoZWFkIjp7Im5hbWUiOiJlcGkifX19XSxbMSwzLCJcXHJyXFxjdXBcXHJzIiwwLHsic3R5bGUiOnsiYm9keSI6eyJuYW1lIjoiZG90dGVkIn0sImhlYWQiOnsibmFtZSI6ImVwaSJ9fX1dLFsyLDMsIlxccnJcXGN1cFxccnMiLDIseyJzdHlsZSI6eyJib2R5Ijp7Im5hbWUiOiJkb3R0ZWQifSwiaGVhZCI6eyJuYW1lIjoiZXBpIn19fV0sWzMsMiwiXFxwcmVjZXEiLDIseyJzdHlsZSI6eyJib2R5Ijp7Im5hbWUiOiJub25lIn0sImhlYWQiOnsibmFtZSI6ImVwaSJ9fX1dLFsxLDAsIlxccHJlY2VxIiwyLHsic3R5bGUiOnsiYm9keSI6eyJuYW1lIjoibm9uZSJ9LCJoZWFkIjp7Im5hbWUiOiJlcGkifX19XV0=
\[\begin{tikzcd}
	a & b \\
	c & d
	\arrow["\rr\cup\rs"', two heads, from=1-1, to=2-1]
	\arrow["\rr\cup\rs"', two heads, from=1-1, to=1-2]
	\arrow["\rr\cup\rs", dotted, two heads, from=1-2, to=2-2]
	\arrow["\rr\cup\rs"', dotted, two heads, from=2-1, to=2-2]
	\arrow["\preceq"', no body, two heads, from=2-2, to=2-1]
	\arrow["\preceq"', no body, two heads, from=1-2, to=1-1]
\end{tikzcd}\]
 \end{lemma}
\begin{proof}
  By induction on the size of the left edge of the diagram. 

  \begin{description}
  \item[Base case] This diagram clearly holds by letting the bottom edge be the same as the top edge:

% https://q.uiver.app/?q=WzAsNCxbMCwwLCJhIl0sWzEsMCwiYiJdLFswLDEsImMiXSxbMSwxLCJkIl0sWzAsMiwiXFxlcXVpdiIsMSx7InN0eWxlIjp7ImJvZHkiOnsibmFtZSI6Im5vbmUifSwiaGVhZCI6eyJuYW1lIjoibm9uZSJ9fX1dLFswLDEsIlxccnJcXGN1cFxccnMiLDIseyJzdHlsZSI6eyJoZWFkIjp7Im5hbWUiOiJlcGkifX19XSxbMSwzLCJcXGVxdWl2IiwxLHsic3R5bGUiOnsiYm9keSI6eyJuYW1lIjoibm9uZSJ9LCJoZWFkIjp7Im5hbWUiOiJub25lIn19fV0sWzIsMywiXFxyclxcY3VwXFxycyIsMix7InN0eWxlIjp7ImJvZHkiOnsibmFtZSI6ImRvdHRlZCJ9LCJoZWFkIjp7Im5hbWUiOiJlcGkifX19XSxbMywyLCJcXHByZWNlcSIsMix7InN0eWxlIjp7ImJvZHkiOnsibmFtZSI6Im5vbmUifSwiaGVhZCI6eyJuYW1lIjoiZXBpIn19fV0sWzEsMCwiXFxwcmVjZXEiLDIseyJzdHlsZSI6eyJib2R5Ijp7Im5hbWUiOiJub25lIn0sImhlYWQiOnsibmFtZSI6ImVwaSJ9fX1dXQ==
\[\begin{tikzcd}
	a & b \\
	c & d
	\arrow["\equiv"{description}, draw=none, from=1-1, to=2-1]
	\arrow["\rr\cup\rs"', two heads, from=1-1, to=1-2]
	\arrow["\equiv"{description}, draw=none, from=1-2, to=2-2]
	\arrow["\rr\cup\rs"', dotted, two heads, from=2-1, to=2-2]
	\arrow["\preceq"', no body, two heads, from=2-2, to=2-1]
	\arrow["\preceq"', no body, two heads, from=1-2, to=1-1]
\end{tikzcd}\]

    and it implies this diagram:

% https://q.uiver.app/?q=WzAsNCxbMCwwLCJhIl0sWzEsMCwiYiJdLFswLDEsImMiXSxbMSwxLCJkIl0sWzAsMiwiKFxccnJcXGN1cFxccnMpXjAiLDJdLFswLDEsIlxccnJcXGN1cFxccnMiLDIseyJzdHlsZSI6eyJoZWFkIjp7Im5hbWUiOiJlcGkifX19XSxbMSwzLCJcXHJyXFxjdXBcXHJzIiwwLHsic3R5bGUiOnsiYm9keSI6eyJuYW1lIjoiZG90dGVkIn0sImhlYWQiOnsibmFtZSI6ImVwaSJ9fX1dLFsyLDMsIlxccnJcXGN1cFxccnMiLDIseyJzdHlsZSI6eyJib2R5Ijp7Im5hbWUiOiJkb3R0ZWQifSwiaGVhZCI6eyJuYW1lIjoiZXBpIn19fV0sWzMsMiwiXFxwcmVjZXEiLDIseyJzdHlsZSI6eyJib2R5Ijp7Im5hbWUiOiJub25lIn0sImhlYWQiOnsibmFtZSI6ImVwaSJ9fX1dLFsxLDAsIlxccHJlY2VxIiwyLHsic3R5bGUiOnsiYm9keSI6eyJuYW1lIjoibm9uZSJ9LCJoZWFkIjp7Im5hbWUiOiJlcGkifX19XV0=
\[\begin{tikzcd}
	a & b \\
	c & d
	\arrow["{(\rr\cup\rs)^0}"', from=1-1, to=2-1]
	\arrow["\rr\cup\rs"', two heads, from=1-1, to=1-2]
	\arrow["\rr\cup\rs", dotted, two heads, from=1-2, to=2-2]
	\arrow["\rr\cup\rs"', dotted, two heads, from=2-1, to=2-2]
	\arrow["\preceq"', no body, two heads, from=2-2, to=2-1]
	\arrow["\preceq"', no body, two heads, from=1-2, to=1-1]
\end{tikzcd}\]
    
	\item[Inductive case] Assume the diagram holds for left edges of all sizes up to $n$. Then this diagram:

% https://q.uiver.app/?q=WzAsNCxbMCwwLCJhIl0sWzEsMCwiYiJdLFswLDEsImMiXSxbMSwxLCJkIl0sWzAsMiwiKFxccnJcXGN1cFxccnMpXntuKzF9IiwyXSxbMCwxLCJcXHJyXFxjdXBcXHJzIiwyLHsic3R5bGUiOnsiaGVhZCI6eyJuYW1lIjoiZXBpIn19fV0sWzEsMywiXFxyclxcY3VwXFxycyIsMCx7InN0eWxlIjp7ImJvZHkiOnsibmFtZSI6ImRvdHRlZCJ9LCJoZWFkIjp7Im5hbWUiOiJlcGkifX19XSxbMiwzLCJcXHJyXFxjdXBcXHJzIiwyLHsic3R5bGUiOnsiYm9keSI6eyJuYW1lIjoiZG90dGVkIn0sImhlYWQiOnsibmFtZSI6ImVwaSJ9fX1dLFszLDIsIlxccHJlY2VxIiwyLHsic3R5bGUiOnsiYm9keSI6eyJuYW1lIjoibm9uZSJ9LCJoZWFkIjp7Im5hbWUiOiJlcGkifX19XSxbMSwwLCJcXHByZWNlcSIsMix7InN0eWxlIjp7ImJvZHkiOnsibmFtZSI6Im5vbmUifSwiaGVhZCI6eyJuYW1lIjoiZXBpIn19fV1d
\[\begin{tikzcd}
	a & b \\
	c & d
	\arrow["{(\rr\cup\rs)^{n+1}}"', from=1-1, to=2-1]
	\arrow["\rr\cup\rs"', two heads, from=1-1, to=1-2]
	\arrow["\rr\cup\rs", dotted, two heads, from=1-2, to=2-2]
	\arrow["\rr\cup\rs"', dotted, two heads, from=2-1, to=2-2]
	\arrow["\preceq"', no body, two heads, from=2-2, to=2-1]
	\arrow["\preceq"', no body, two heads, from=1-2, to=1-1]
\end{tikzcd}\]

    \noindent
    follows from this diagram:

% https://q.uiver.app/?q=WzAsNixbMCwwLCJhIl0sWzEsMCwiYiJdLFswLDEsImMiXSxbMSwxLCJkIl0sWzAsMiwiZSJdLFsxLDIsImYiXSxbMCwyLCIoXFxyclxcY3VwXFxycylebiIsMl0sWzAsMSwiXFxyclxcY3VwXFxycyIsMix7InN0eWxlIjp7ImhlYWQiOnsibmFtZSI6ImVwaSJ9fX1dLFsxLDMsIlxccnJcXGN1cFxccnMiLDAseyJzdHlsZSI6eyJib2R5Ijp7Im5hbWUiOiJkb3R0ZWQifSwiaGVhZCI6eyJuYW1lIjoiZXBpIn19fV0sWzIsMywiXFxyclxcY3VwXFxycyIsMix7InN0eWxlIjp7ImJvZHkiOnsibmFtZSI6ImRvdHRlZCJ9LCJoZWFkIjp7Im5hbWUiOiJlcGkifX19XSxbMywyLCJcXHByZWNlcSIsMix7InN0eWxlIjp7ImJvZHkiOnsibmFtZSI6Im5vbmUifSwiaGVhZCI6eyJuYW1lIjoiZXBpIn19fV0sWzEsMCwiXFxwcmVjZXEiLDIseyJzdHlsZSI6eyJib2R5Ijp7Im5hbWUiOiJub25lIn0sImhlYWQiOnsibmFtZSI6ImVwaSJ9fX1dLFsyLDQsIlxccnJcXGN1cFxccnMiLDJdLFszLDUsIlxccnJcXGN1cFxccnMiLDAseyJzdHlsZSI6eyJib2R5Ijp7Im5hbWUiOiJkb3R0ZWQifSwiaGVhZCI6eyJuYW1lIjoiZXBpIn19fV0sWzQsNSwiXFxyclxcY3VwXFxycyIsMix7InN0eWxlIjp7ImJvZHkiOnsibmFtZSI6ImRvdHRlZCJ9LCJoZWFkIjp7Im5hbWUiOiJlcGkifX19XSxbNSw0LCJcXHByZWNlcSIsMix7InN0eWxlIjp7ImJvZHkiOnsibmFtZSI6Im5vbmUifSwiaGVhZCI6eyJuYW1lIjoiZXBpIn19fV0sWzEyLDEzLCJcXGxhYmVsY3JlZntsZW06c2tld3Jvd30iLDEseyJzaG9ydGVuIjp7InNvdXJjZSI6MjAsInRhcmdldCI6MjB9LCJzdHlsZSI6eyJib2R5Ijp7Im5hbWUiOiJub25lIn0sImhlYWQiOnsibmFtZSI6Im5vbmUifX19XV0=
\[\begin{tikzcd}
	a & b \\
	c & d \\
	e & f
	\arrow["{(\rr\cup\rs)^n}"', from=1-1, to=2-1]
	\arrow["\rr\cup\rs"', two heads, from=1-1, to=1-2]
	\arrow["\rr\cup\rs", dotted, two heads, from=1-2, to=2-2]
	\arrow["\rr\cup\rs"', dotted, two heads, from=2-1, to=2-2]
	\arrow["\preceq"', no body, two heads, from=2-2, to=2-1]
	\arrow["\preceq"', no body, two heads, from=1-2, to=1-1]
	\arrow[""{name=0, anchor=center, inner sep=0}, "\rr\cup\rs"', from=2-1, to=3-1]
	\arrow[""{name=1, anchor=center, inner sep=0}, "\rr\cup\rs", dotted, two heads, from=2-2, to=3-2]
	\arrow["\rr\cup\rs"', dotted, two heads, from=3-1, to=3-2]
	\arrow["\preceq"', no body, two heads, from=3-2, to=3-1]
	\arrow["{\labelcref{lem:skewrow}}"{description}, draw=none, from=0, to=1]
\end{tikzcd}\]


\noindent
  where the top half is the inductive hypothesis and the bottom half follows from \cref{lem:skewrow}.
\end{description}
\end{proof}


\begin{lemma}\label{lem:skew-union}
  Let relation $\rr$ be monotonic in quasi order $\preceq$ and skew confluent using $\preceq$ and $\rr_\preceq^\leftarrow$;
  similarly let $\rs$ be monotonic in the same quasi order $\preceq$ and skew confluent using $\preceq$ and $\rs_\preceq^\leftarrow$.
  If $\rr$ commutes with $\rs$, then $\rt=\rr\cup\rs$ is monotonic in $\preceq$ and skew confluent using $\preceq$ and $\rt_\preceq^\leftarrow$.
\end{lemma}
\begin{proof}

\end{proof}
